\documentclass[11pt]{amsart}

\usepackage[T1]{fontenc}
\usepackage[utf8]{inputenc}
\usepackage{lmodern}
\usepackage{microtype}
\usepackage{amsmath,amssymb,amsthm,mathtools,mathrsfs}
\usepackage{enumitem}
\usepackage{aliascnt}
\usepackage[hidelinks]{hyperref}
\usepackage[nameinlink,noabbrev]{cleveref}

\allowdisplaybreaks
\numberwithin{equation}{section}

\newtheorem{theorem}{Theorem}[section]

\newaliascnt{proposition}{theorem}
\newtheorem{proposition}[proposition]{Proposition}
\aliascntresetthe{proposition}

\newaliascnt{lemma}{theorem}
\newtheorem{lemma}[lemma]{Lemma}
\aliascntresetthe{lemma}

\newaliascnt{corollary}{theorem}
\newtheorem{corollary}[corollary]{Corollary}
\aliascntresetthe{corollary}

\newaliascnt{claim}{theorem}
\newtheorem{claim}[claim]{Claim}
\aliascntresetthe{claim}

\theoremstyle{definition}
\newaliascnt{definition}{theorem}
\newtheorem{definition}[definition]{Definition}
\aliascntresetthe{definition}

\theoremstyle{remark}
\newaliascnt{remark}{theorem}
\newtheorem{remark}[remark]{Remark}
\aliascntresetthe{remark}

\crefname{theorem}{theorem}{theorems}
\Crefname{theorem}{Theorem}{Theorems}
\crefname{proposition}{proposition}{propositions}
\Crefname{proposition}{Proposition}{Propositions}
\crefname{lemma}{lemma}{lemmas}
\Crefname{lemma}{Lemma}{Lemmas}
\crefname{corollary}{corollary}{corollaries}
\Crefname{corollary}{Corollary}{Corollaries}
\crefname{claim}{claim}{claims}
\Crefname{claim}{Claim}{Claims}
\crefname{definition}{definition}{definitions}
\Crefname{definition}{Definition}{Definitions}
\crefname{remark}{remark}{remarks}
\Crefname{remark}{Remark}{Remarks}

\newcommand{\R}{\mathbb R}
\newcommand{\Q}{\mathbb Q}
\newcommand{\Z}{\mathbb Z}
\newcommand{\C}{\mathbb C}
\newcommand{\N}{\mathbb N}
\newcommand{\Lio}{\mathscr L}
\newcommand{\Span}{\operatorname{span}}
\newcommand{\ord}{\operatorname{ord}}
\newcommand{\dist}{\operatorname{dist}}
\newcommand{\diam}{\operatorname{diam}}
\newcommand{\Int}{\operatorname{int}}
\newcommand{\rank}{\operatorname{rank}}
\newcommand{\norm}[1]{\left\lVert #1\right\rVert}
\newcommand{\abs}[1]{\left\lvert #1\right\rvert}
\newcommand{\set}[1]{\left\{#1\right\}}
\newcommand{\eps}{\varepsilon}

\title[Mahler's Problem on Liouville Numbers]
{Mahler's Problem on Liouville Numbers}

\author{Diego Marques}
\address{Departamento de Matem\'atica, Universidade de Bras\'ilia, Bras\'ilia, 70910-900, Brazil}
\email{diego@mat.unb.br}
\subjclass[2020]{Primary 11J82; Secondary 11J81, 26E05, 30D20}
\keywords{Liouville numbers, Mahler's problem, Maillet's property,
real-analytic functions, two-height approximation, arithmetic rigidity,
irrationality exponent}
\date{}

\hypersetup{
  pdftitle={Mahler's Problem on Liouville Numbers},
  pdfauthor={Diego Marques},
  pdfsubject={Mahler's Problem on Liouville Numbers},
  pdfkeywords={Liouville numbers, Mahler's problem, Maillet's property,
real-analytic functions, two-height approximation, arithmetic rigidity,
irrationality exponent}
}

\begin{document}

\begin{abstract}
A classical theorem of Maillet states that every nonconstant rational
function with rational coefficients maps Liouville numbers to Liouville numbers.
In 1984, Mahler asked whether there exists a transcendental entire
function with the same property. We resolve this question in the
negative. More generally, we prove that every real-analytic function
on an interval with this property is the restriction of a rational
function in $\R(x)$. The local result is quantitative: for every nonrational
real-analytic function, every nonempty open subinterval of its
domain contains a Liouville number whose image has irrationality
exponent at most $100$. The proof rests on a two-height counting estimate that treats
source and target denominators independently. An adaptive
determinant argument, uniform Wronskian sublevel bounds, and
Farey separation establish the estimate; a nested-interval
construction yields the quantitative result.
\end{abstract}


\maketitle

\section{Introduction}

A real number $\xi$ is called a \emph{Liouville number} if, for every
integer $N\geq1$, there exist infinitely many pairs
$(p,q)\in\Z\times\Z_{>0}$, with $q\geq2$, such that
\begin{equation}\label{eq:intro-liouville}
0<\abs{\xi-\frac pq}<q^{-N}.
\end{equation}
We denote the set of Liouville numbers by $\Lio$. Since Liouville introduced these numbers
\cite{Liouville1844,Liouville1851}, their exceptional
rational approximation properties have made them a natural setting for
the study of arithmetic properties of functions; see, for instance,
\cite{Bugeaud2004,FischlerRivoal2024,Waldschmidt2003}.

A classical result of Maillet \cite{Maillet1906} states that every
nonconstant rational function $R\in\Q(x)$ maps $\Lio$ into itself. We say that a
function has \emph{Maillet's property} if it preserves every Liouville
number in its domain.

In one of his last papers, published in 1984, Mahler
\cite{Mahler1984} formulated several problems at the interface of
transcendence and Diophantine approximation. Motivated by Maillet's
theorem, he asked:

\begin{quote}
\emph{Which analytic functions $f(z)$ have the property that if $\xi$ is
any Liouville number, then so is $f(\xi)$? In particular, are there entire
transcendental functions with this property?}
\end{quote}

The final sentence isolates the central transcendental problem: whether Maillet's arithmetic invariance can persist for a genuinely transcendental entire function. Subsequently recorded by Waldschmidt in his collection of open Diophantine problems \cite[Question~3.27, p.~281]{Waldschmidt2004} and by Bugeaud in his monograph \cite[Problem~44, p.~215]{Bugeaud2004}, this existence problem has remained open for over four decades.

We answer Mahler's question in the negative.

\begin{theorem}\label{thm:intro-mahler}
Let $F:\C\to\C$ be entire. If
\[
F(\Lio)\subseteq\Lio,
\]
then $F\in\R[z]$. In particular, no transcendental entire function has
Maillet's property.
\end{theorem}

%Theorem~\ref{thm:intro-mahler} is a rigidity theorem rather than an
%interpolation result. We impose no assumption that $F(\Q)\subseteq\Q$,
%no restriction on the denominators or heights of values of $F$ at
%rational points, and no growth condition on $F$.

The requirement that \emph{every} Liouville number be preserved
is essential in Theorem~\ref{thm:intro-mahler}. Classical arithmetic
interpolation allows considerable flexibility on countable sets
\cite{Stackel1895,Stackel1902,Weierstrass1923}: transcendental entire
functions can even preserve every algebraic number field
\cite{VanderPoorten1968}; see also
\cite{BoxallJones2015,Lombardo2017,Surroca2006}.
Preservation of Liouville numbers on large sets is likewise
compatible with much greater flexibility. Alnia\c{c}{\i}k and Saias
\cite{AlniacikSaias1994} showed that a continuous, nowhere locally
constant real function maps a dense $G_\delta$ subset of $\Lio$
in its domain into $\Lio$, and transcendental entire functions
can preserve distinguished subclasses of $\Lio$
\cite{MarquesMoreira2015,MarquesSchleischitz2016}.

The regularity boundary is sharp within the classical
smooth--analytic hierarchy. The flexibility established at finite
orders of smoothness \cite{LelisMoreiraSilva2026} persists in the
$C^\infty$ category: there exist transcendental $C^\infty$
diffeomorphisms of $\mathbb R$, arbitrarily close to the identity
in the compact-open $C^\infty$ topology, such that every derivative
of every positive iterate has Maillet's property
\cite{MarquesSmooth2026}. By contrast,
Corollary~\ref{cor:intro-rigidity} shows that real analyticity
forces rational rigidity.

Theorem~\ref{thm:intro-mahler} follows from a local quantitative
result. Recall that the irrationality exponent of
$y\in\mathbb R$ is
\begin{equation}\label{eq:intro-mu}
\mu(y):=
\sup\left\{
\lambda>0:
\begin{aligned}
&0<\abs{y-a/b}<b^{-\lambda}\\
&\text{for infinitely many }(a,b)\in\Z\times\Z_{>0}
\end{aligned}
\right\}.
\end{equation}
Thus
\[
y\in\Lio\quad\Longleftrightarrow\quad\mu(y)=+\infty.
\]

\begin{theorem}\label{thm:intro-local}
There exists an absolute constant $\tau>2$ with the following property.
Let $U\subseteq\mathbb R$ be an open interval and let
$f:U\to\mathbb R$ be real-analytic. Suppose that $f$ is not the
restriction to $U$ of a rational function in $\mathbb R(x)$ without
poles on $U$. Then every nonempty open subinterval $V\subseteq U$
contains a point $\xi\in\Lio$ such that
\[
\mu(f(\xi))\leq\tau.
\]
In fact, $\xi$ may be chosen so that
\[
\left|f(\xi)-\frac{a}{b}\right|>b^{-\tau}
\]
for every $a\in\mathbb Z$ and every sufficiently large integer $b$.
One may take $\tau=100$.
\end{theorem}

The particular choice $\tau=100$ is not intrinsic and has not been
optimized. The substantive point is the existence of an absolute
finite bound, independent of the function. Taking the contrapositive of
Theorem~\ref{thm:intro-local} gives the corresponding local rigidity
statement.

\begin{corollary}\label{cor:intro-rigidity}
Let $f:U\to\R$ be real-analytic on an open interval. If
\[
f(\Lio\cap U)\subseteq\Lio,
\]
then there exists a rational function $R\in\R(x)$, with no pole in $U$,
such that
\[
f(x)=R(x)\qquad(x\in U).
\]
\end{corollary}

The rational-rigidity conclusion raises a natural arithmetic question:
must the resulting rational function be defined over $\Q$?  The answer
displays a sharp distinction between degree one and higher degree.

\begin{theorem}\label{thm:intro-mobius}
Let
\[
M(x)=\frac{ax+b}{cx+d}\in\R(x),
\qquad ad-bc\neq0.
\]
Then $M$ has Maillet's property if and only if $M\in\Q(x)$,
equivalently, if and only if
\[
[a:b:c:d]\in\mathbb P^3(\Q).
\]
Moreover, if $M\notin\Q(x)$, then every nonempty open interval
contained in the domain of $M$ contains a point $\xi\in\Lio$ such that
$M(\xi)\notin\Lio$.
\end{theorem}

\begin{theorem}\label{thm:intro-nonlinear}
There exist continuum many positive transcendental real numbers $\sigma$
with the following property: for every integer $m\geq2$ and every
$\xi\in\Lio$,
\[
\sigma\xi^m\in\Lio.
\]
\end{theorem}

Thus coefficient rigidity holds for rational maps of degree one but
fails in every degree at least two. In particular, the conclusion
$F\in\R[z]$ in Theorem~\ref{thm:intro-mahler} cannot be strengthened
to $F\in\overline{\Q}[z]$. We do not attempt here to classify all rational functions in $\R(x)$ with Maillet's property.

The passage to the entire case is immediate. An entire function
preserving $\Lio$ is real-valued on $\mathbb R$ by density and
continuity. Corollary~\ref{cor:intro-rigidity} makes its restriction
rational; analytic continuation and the absence of poles then
force it to be a polynomial.\\

\noindent
\textbf{The two-height mechanism.}
The local theorem requires arbitrarily good rational approximation
in the source and exclusion of very good rational approximation
in the image at every sufficiently large denominator. We construct
such a point by counting source fractions $p/q$, with $q\asymp Q$,
whose images are close to target fractions $a/b$, keeping the two
denominator scales separate.

Small target denominators are handled by elementary counting.
For larger denominators, write $b\asymp Q^u$ and adapt the source
degree to $u$, while keeping the target degree equal to one.
A determinant comparison places the rational pairs in each short
source cell on a common relation
\[
A(x)+B(x)y=0.
\]
Such a relation cannot hold identically on the graph of a
nonrational $f$; this is why the argument also applies to
nonrational algebraic functions. After localization away from the
relevant Wronskian zeros, sublevel estimates uniform in the
normalized coefficients control both the length and the number
of components of the set where $A+Bf$ is small. Farey separation
then gives a bounded number of source fractions per cell, while
the choice of degree keeps the number of cells subquadratic
in $Q$.

The resulting estimate has a quadratic term controlled by the
lower target cutoff and a subquadratic remainder uniform in that
cutoff. The constant multiplying the quadratic term is independent
of the source subinterval. These uniformities allow safe rational
centers to be chosen repeatedly on shrinking intervals.

We then choose nested intervals near these centers. The source
approximation orders tend to infinity, and target avoidance is
imposed with a margin that persists throughout each interval.
Consecutive target-denominator ranges cover all sufficiently
large denominators. The limiting point is therefore Liouville,
while its image has irrationality exponent at most $100$.
Liouville numbers enter only in this final construction; the
counting estimate concerns rational approximation to the graph
itself.

The determinant argument lies in the tradition of Bombieri--Pila and
Pila \cite{BombieriPila1989,Pila1991,Pila2005}; related developments
include \cite{BeresnevichDickinsonVelani2007,ComteYomdin2018,
Habegger2018,Huang2015,Huxley1994,KleinbockMargulis1998,
PilaWilkie2006,VaughanVelani2006}. Separate coordinate denominators and Wronskian arguments already
occur in Pila's work. 

The results of Section~\ref{sec:arithmetic-boundary} are independent of the analytic counting
argument. The degree-one classification uses a Diophantine
dichotomy for the coefficient vector and rational reconstruction
from three points. The higher-degree examples combine a
continued-fraction construction, Petruska's theorem on strong
Liouville numbers, and a gap principle for $U_m$-numbers.

\section{Rational counting preliminaries}
\label{sec:rational-counting}

\noindent
\textbf{Notation and conventions.} For a finite union of intervals $E\subseteq\R$, we denote its Lebesgue
measure by $|E|$. We write $J\Subset U$ if $\overline J$ is a compact
subset of $U$.

The symbols $O(\cdot)$, $\ll$, and $\asymp$ have their usual meanings.
Unless otherwise stated, implied constants may depend on quantities
fixed before the estimate, but not on parameters subsequently allowed
to vary or tend to infinity. Additional dependence is indicated by
subscripts; thus $X\ll_{A,J}Y$ means that
$X\leq C(A,J)Y$ for some $C(A,J)>0$. Generic constants
$C,C_1,C_2,\ldots$ may vary from one occurrence to the next.

Subscripts or superscripts on named constants are mnemonic rather than
indicating parameter dependence; for example, \textup{sub} and
\textup{low} refer to sublevel and low-target-height estimates,
respectively.

With these conventions in place, we record the two rational-counting
estimates needed in Section~\ref{sec:scarcity}. They provide complementary
upper and lower bounds for reduced rational points of comparable height.
The first is based on the separation of reduced fractions; the absence
of a factor $Q$ in the component term will be essential in Section~\ref{sec:scarcity}.

\begin{lemma}[Reduced fractions in a finite union of intervals]
\label{lem:rational-union}
Let $E\subseteq\R$ be a union of at most $R$ intervals, where
$R\in\Z_{>0}$, and let $Q\in\Z_{>0}$. Reduced fractions are counted as rational
numbers, so their positive-denominator representations are unique. Then
\begin{equation}\label{eq:rational-union}
\#\set{\frac pq\in E:\gcd(p,q)=1,\ Q\leq q<2Q}
\leq4Q^2|E|+R.
\end{equation}
\end{lemma}

\begin{proof}
Write $E$ as the disjoint union of its nonempty connected components
$I_1,\ldots,I_s$, where $s\leq R$. If $p/q\neq p'/q'$ are reduced
fractions with positive denominators $q,q'<2Q$, then
\[
\abs{\frac pq-\frac{p'}{q'}}
=
\frac{|pq'-p'q|}{qq'}
\geq\frac1{qq'}>\frac1{4Q^2}.
\]
Consequently, an interval of length $\ell$ contains at most
$4Q^2\ell+1$ such fractions. Applying this to each $I_\nu$ and
summing proves \eqref{eq:rational-union}.
\end{proof}

The complementary estimate supplies a positive proportion of reduced
source fractions in every fixed interval.

\begin{lemma}[Reduced fractions in a denominator block]
\label{lem:farey-lower-en}
There exists an absolute constant $c_{\mathrm F}>0$ with the following
property. For every nondegenerate compact interval $J\subseteq\R$,
there exists $Q_0(J)\in\Z_{>0}$ such that
\begin{equation}\label{eq:farey-lower-en}
\#\set{\frac pq\in J:\gcd(p,q)=1,\ Q\leq q<2Q}
\geq c_{\mathrm F}|J|Q^2
\end{equation}
for every integer $Q\geq Q_0(J)$. One may take
$c_{\mathrm F}=1/4$.
\end{lemma}

\begin{proof}
Write $J=[\alpha,\beta]$ and, for $q\ge1$, let
\[
N_q(J):=\#\{p\in\Z:\alpha q\le p\le\beta q,\ \gcd(p,q)=1\}.
\]
By M\"obius inversion,
\[
N_q(J)=|J|\varphi(q)+O(\tau_0(q)),
\]
where $\tau_0(q)$ denotes the number of positive
divisors of $q$, and the implied constant is absolute. Hence, using the classical estimates \cite{Apostol1976},
\[
\sum_{n\le X}\varphi(n)=\frac{3}{\pi^2}X^2+O(X\log X),
\qquad
\sum_{n\le X}\tau_0(n)=O(X\log X),
\]
we obtain
\[
\sum_{Q\le q<2Q}N_q(J)
=
\frac{9}{\pi^2}|J|Q^2
+O\bigl((1+|J|)Q\log Q\bigr).
\]
Since $J$ is fixed and nondegenerate, the error term is
$o(|J|Q^2)$. As $9/\pi^2>1/4$, the result follows for all
sufficiently large $Q$.
\end{proof}


\section{Wronskian localization and uniform sublevel bounds}
\label{sec:wronskians}

In Section~\ref{sec:scarcity}, the determinant argument produces auxiliary
polynomials $A(X)+B(X)Y$ whose coefficients depend on the
rational points being counted. We therefore need sublevel
estimates for $A(x)+B(x)f(x)$ that are uniform over all normalized
coefficient vectors. Nonrationality of $f$ provides the required
Wronskian nondegeneracy after localization.

Let $\phi_1,\ldots,\phi_N$ be functions of class $C^{N-1}$ on an
interval. Their \textit{Wronskian} is
\begin{equation}\label{eq:wronskian-definition-en}
W(\phi_1,\ldots,\phi_N)(x)
:=
\det\bigl(\phi_j^{(k-1)}(x)\bigr)_{1\leq k,j\leq N}.
\end{equation}

\subsection{Rationality and linear independence}

We first recall the classical Wronskian criterion for analytic
functions. The analyticity assumption is essential here: the
corresponding statement fails, in general, for smooth functions.

\begin{lemma}[Analytic independence and Wronskians]
\label{lem:analytic-wronskian-en}
Let $U\subseteq\R$ be a connected open interval. If
$\phi_1,\ldots,\phi_N$ are real-analytic and linearly independent on
$U$, then
\[
W(\phi_1,\ldots,\phi_N)\not\equiv0
\qquad\text{on }U.
\]
\end{lemma}

This is the classical analytic Wronskian criterion, going back at least
to B\^ocher \cite{Bocher1900}; see also \cite{BostanDumas2010} for a modern
treatment.

For the remainder of the local argument, fix a real-analytic function
$f:U\to\R$ which is not the restriction to $U$ of a rational function
in $\R(x)$. For $d\geq0$, put
\begin{equation}\label{eq:N-d-en}
N_d:=2(d+1)
\end{equation}
and consider the indexed family
\begin{equation}\label{eq:rational-family-en}
\mathcal F_d(f)
:=
\bigl(x^i f(x)^j\bigr)_{\substack{0\leq i\leq d\\ j\in\{0,1\}}}.
\end{equation}

\begin{lemma}[Rationality and linear dependence]
\label{lem:rational-independence-en}
Let $f:U\to\R$ be real-analytic. The following are equivalent.
\begin{enumerate}[label=\textup{(\roman*)},leftmargin=2.7em]
\item $f$ is the restriction to $U$ of a rational function in
$\R(x)$ with no pole in $U$;
\item the family $\mathcal F_d(f)$ is linearly dependent over $\R$
for some $d\geq0$.
\end{enumerate}
Consequently, under the standing hypothesis on $f$, the family
$\mathcal F_d(f)$ is linearly independent for every $d\geq0$.
\end{lemma}

\begin{proof}
Suppose first that $f=P/Q$ on $U$, where $P,Q\in\R[X]$, $Q$ has no
zero in $U$, and $Q\neq0$. If
$d\geq\max\{\deg P,\deg Q\}$, then
\[
Q(x)f(x)-P(x)=0
\qquad(x\in U),
\]
which is a nontrivial linear relation among the members of
$\mathcal F_d(f)$.

Conversely, suppose that $\mathcal F_d(f)$ is linearly dependent. Then
there are polynomials $A,B\in\R[X]$, not both zero and of degree at
most $d$, such that
\begin{equation}\label{eq:A-B-f-relation-en}
A(x)+B(x)f(x)=0
\qquad(x\in U).
\end{equation}
If $B=0$, then $A$ vanishes on the open interval $U$, and hence
$A=0$, a contradiction. Thus $B\neq0$.

Let $D=\gcd(A,B)$ in $\R[X]$ and write $A=DA_0$, $B=DB_0$, with
$A_0$ and $B_0$ coprime. On the nonempty open set
$U\setminus\{D=0\}$, equation \eqref{eq:A-B-f-relation-en} gives
$A_0+B_0f=0$. Since $D$ has only finitely many zeros, continuity extends this identity to all of $U$:
\[
A_0(x)+B_0(x)f(x)=0
\qquad(x\in U).
\]
If $B_0(x_0)=0$ at some $x_0\in U$, then the last identity also gives
$A_0(x_0)=0$, contrary to the coprimality of $A_0$ and $B_0$.
Therefore $B_0$ has no zero in $U$, and
\[
f(x)=-\frac{A_0(x)}{B_0(x)}
\qquad(x\in U).
\]
This proves \textup{(i)}.
\end{proof}

Order the pairs
\[
(i,j)\in\{0,\ldots,d\}\times\{0,1\}
\]
lexicographically, and denote the corresponding ordered family of
functions $x^if(x)^j$ by
\[
\varphi_{d,1},\ldots,\varphi_{d,N_d}.
\]
Define
\begin{equation}\label{eq:W-d-en}
W_d(x)
:=
W(\varphi_{d,1},\ldots,\varphi_{d,N_d})(x).
\end{equation}
The ordering affects $W_d$ at most by a sign and therefore has no
bearing on its zero set.

\begin{corollary}\label{cor:wronskian-finite-en}
For every $d\geq0$, the function $W_d$ is real-analytic and not
identically zero on $U$. Consequently, if $I\Subset U$ is a compact
interval, then
\[
\set{x\in I:W_d(x)=0}
\]
is finite.
\end{corollary}

\begin{proof}
The functions in $\mathcal F_d(f)$ are real-analytic and, by
Lemma~\ref{lem:rational-independence-en}, linearly independent.
Lemma~\ref{lem:analytic-wronskian-en} therefore gives
$W_d\not\equiv0$ on $U$. The function $W_d$ is real-analytic because
it is a determinant of derivatives of real-analytic functions. If it
had infinitely many zeros in $I$, compactness would give an
accumulation point in $U$, and the identity theorem would force
$W_d\equiv0$ on the connected interval $U$, a contradiction.
\end{proof}

\subsection{One-dimensional sublevel estimates}

We next record the elementary one-dimensional estimate underlying the
uniform argument. It expresses the principle that a function with a
derivative uniformly separated from zero cannot remain small on a
long set.

\begin{lemma}\label{lem:one-derivative-en}
Let $k\geq1$, let $I\subseteq\R$ be a nondegenerate compact interval,
and let $g\in C^k(I)$. Suppose that
\begin{equation}\label{eq:kth-lower-en}
\abs{g^{(k)}(x)}\geq\lambda>0
\qquad(x\in I).
\end{equation}
Then, for every $\eps>0$, the set
\[
E_g(\eps):=\set{x\in I:\abs{g(x)}\leq\eps}
\]
has at most $2k+1$ connected components and satisfies
\begin{equation}\label{eq:one-derivative-length-en}
|E_g(\eps)|
\leq
2k(2k+1)
\left(\frac{\eps}{\lambda}\right)^{1/k}.
\end{equation}
\end{lemma}

\begin{proof}
Since $g^{(k)}$ is continuous and nowhere zero on $I$, it has constant
sign. By Rolle's theorem, for every $c\in\R$ the equation $g(x)=c$
has at most $k$ solutions. Hence the relative boundary of
\[
E_g(\varepsilon)=\{x\in I:|g(x)|\leq\varepsilon\}
\]
has at most $2k$ points, and $E_g(\varepsilon)$ has at most
$2k+1$ connected components.

Let $[a,b]$ be a nondegenerate component of $E_g(\varepsilon)$, and put
\[
h:=\frac{b-a}{k},
\qquad
x_j:=a+jh
\quad(0\leq j\leq k).
\]
By the mean-value theorem for finite differences, for some
$\theta\in(a,b)$,
\[
h^k g^{(k)}(\theta)
=
\sum_{j=0}^k(-1)^{k-j}\binom{k}{j}g(x_j).
\]
Since $|g(x_j)|\leq\varepsilon$,
\[
h^k\lambda
\leq
2^k\varepsilon,
\]
and therefore
\[
b-a=kh
\leq
2k\left(\frac{\varepsilon}{\lambda}\right)^{1/k}.
\]
Summing over at most $2k+1$ nondegenerate components gives
\[
|E_g(\varepsilon)|
\leq
2k(2k+1)
\left(\frac{\varepsilon}{\lambda}\right)^{1/k}.
\]
\end{proof}

We now pass from a lower bound for one prescribed derivative to a
bound that is uniform over all normalized linear combinations of a
fixed finite family. Wronskian nondegeneracy gives a uniform lower bound for at least
one derivative at each point. Compactness of the product of the
coefficient sphere and the source interval makes the resulting
local estimates uniform.

\begin{lemma}[Uniform sublevel bounds under Wronskian nondegeneracy]
\label{lem:uniform-sublevel-en}
Let $J\subseteq\R$ be a nondegenerate compact interval, let $N\geq2$,
and let $\phi_1,\ldots,\phi_N$ be of class $C^{N-1}$ on a neighborhood
of $J$. Assume that
\begin{equation}\label{eq:wronskian-separated-en}
\inf_{x\in J}\abs{W(\phi_1,\ldots,\phi_N)(x)}>0.
\end{equation}
For $c=(c_1,\ldots,c_N)\in\R^N$, write
\[
G_c(x)=\sum_{j=1}^N c_j\phi_j(x).
\]
Then there exist constants
\[
C_{\mathrm{sub}}>0,
\qquad
R_{\mathrm{sub}}\in\Z_{>0},
\qquad
\eps_0\in(0,1)
\]
such that, for every $c\in\R^N$ with $\norm{c}_2=1$ and every
$0<\eps<\eps_0$, the set
\[
\set{x\in J:\abs{G_c(x)}\leq\eps}
\]
is a union of at most $R_{\mathrm{sub}}$ intervals and satisfies
\begin{equation}\label{eq:uniform-sublevel-en}
\left|\set{x\in J:\abs{G_c(x)}\leq\eps}\right|
\leq C_{\mathrm{sub}}\eps^{1/(N-1)}.
\end{equation}
\end{lemma}

\begin{proof}
Consider the jet matrix
\[
\mathcal A(x)
=
\bigl(\phi_j^{(k)}(x)\bigr)_{
\substack{0\leq k\leq N-1\\1\leq j\leq N}}.
\]
Its determinant is the Wronskian in
\eqref{eq:wronskian-separated-en}; hence $\mathcal A(x)$ is invertible
for every $x\in J$. The least singular value of $\mathcal A(x)$
depends continuously on $x$. Since $J$ is compact, there exists
$\sigma>0$ such that
\begin{equation}\label{eq:uniform-singular-value-en}
\norm{\mathcal A(x)c}_2\geq\sigma\norm{c}_2
\qquad(x\in J,\ c\in\R^N).
\end{equation}
For $\norm{c}_2=1$,
\[
\mathcal A(x)c
=
\bigl(G_c(x),G_c'(x),\ldots,G_c^{(N-1)}(x)\bigr)^{\mathsf T},
\]
and therefore
\[
\max_{0\leq k\leq N-1}\abs{G_c^{(k)}(x)}
\geq
\eta,
\qquad
\eta:=\frac{\sigma}{\sqrt N},
\]
for every $(c,x)\in\mathbb S^{N-1}\times J$.

Fix $(c_0,x_0)\in\mathbb S^{N-1}\times J$. There is an index
$k(c_0,x_0)\in\{0,\ldots,N-1\}$ for which
\[
\abs{G_{c_0}^{(k(c_0,x_0))}(x_0)}\geq\eta.
\]
Continuity in $(c,x)$ yields a relatively open neighborhood $U$ of
$c_0$ in $\mathbb S^{N-1}$ and a compact interval $I\subseteq J$,
with $x_0\in\Int_J(I)$, such that
\[
\abs{G_c^{(k(c_0,x_0))}(x)}\geq\frac{\eta}{2}
\qquad(c\in U,\ x\in I).
\]
These sets $U\times\Int_J(I)$ cover the compact space
$\mathbb S^{N-1}\times J$. Choose a finite subcover
\[
U_\nu\times\Int_J(I_\nu),
\qquad1\leq\nu\leq L,
\]
and let $k_\nu\in\{0,\ldots,N-1\}$ denote the corresponding derivative
order. Thus
\begin{equation}\label{eq:local-large-jet-en}
\abs{G_c^{(k_\nu)}(x)}\geq\frac{\eta}{2}
\qquad(c\in U_\nu,\ x\in I_\nu).
\end{equation}
The finite cover, and hence $L$ and all constants derived from it,
depend only on $J$ and the family $\phi_1,\ldots,\phi_N$, and not on
the subsequent choice of $c\in\mathbb S^{N-1}$.

Fix $c\in\mathbb S^{N-1}$ and let
\[
\mathcal N(c):=\set{\nu\in\{1,\ldots,L\}:c\in U_\nu}.
\]
Then the intervals $I_\nu$, $\nu\in\mathcal N(c)$, cover $J$ by their
relative interiors. Put
\[
\eps_0:=\min\left\{\frac12,\frac{\eta}{4}\right\}.
\]
If $0<\eps<\eps_0$ and $k_\nu=0$, then
\eqref{eq:local-large-jet-en} shows that the sublevel set does not meet
$I_\nu$. If $1\leq k_\nu\leq N-1$, Lemma~\ref{lem:one-derivative-en},
with $\lambda=\eta/2$, gives
\[
\left|\set{x\in I_\nu:\abs{G_c(x)}\leq\eps}\right|
\leq
2k_\nu(2k_\nu+1)
\left(\frac{2\eps}{\eta}\right)^{1/k_\nu}.
\]
Since $0<\eps<1$ and $k_\nu\leq N-1$, this is at most
\[
C_*\eps^{1/(N-1)},
\qquad
C_*:=
\max_{1\leq k\leq N-1}
2k(2k+1)\left(\frac{2}{\eta}\right)^{1/k}.
\]
The same lemma shows that each such local sublevel set has at most
$2N-1$ connected components.

Since $\#\mathcal N(c)\leq L$, subadditivity of length now gives
\[
\left|\set{x\in J:\abs{G_c(x)}\leq\eps}\right|
\leq
LC_*\eps^{1/(N-1)}.
\]
Moreover, this set has at most $L(2N-1)$ connected components.
Thus one may take
\[
C_{\mathrm{sub}}:=LC_*,
\qquad
R_{\mathrm{sub}}:=L(2N-1),
\]
and these constants, as well as $\eps_0$, depend only on $J$ and the
family $\phi_1,\ldots,\phi_N$, and not on $c$.
\end{proof}

We specialize this result to the target-linear polynomials that arise
from the determinant argument. For
\[
P(X,Y)=A(X)+B(X)Y
=\sum_{i=0}^d c_{i0}X^i+\sum_{i=0}^d c_{i1}X^iY,
\]
write
\[
\norm{P}_2
:=
\left(\sum_{i=0}^d(c_{i0}^2+c_{i1}^2)\right)^{1/2}.
\]

\begin{corollary}[Uniform rational-relation sublevel estimate]
\label{cor:polynomial-sublevel-en}
Fix $d\geq0$ and a nondegenerate compact interval $J\Subset U$ on which $W_d$ does
not vanish. Put $N=N_d=2(d+1)$. Then there exist constants
\[
C_{d,J}>0,
\qquad
R_{d,J}\in\Z_{>0},
\qquad
\eps_{d,J}>0
\]
such that, for every polynomial
\[
P(X,Y)=A(X)+B(X)Y,
\qquad
\deg A,\deg B\leq d,
\qquad
\norm{P}_2=1,
\]
and every $0<\eps<\eps_{d,J}$, the set
\[
\set{x\in J:\abs{P(x,f(x))}\leq\eps}
\]
is a union of at most $R_{d,J}$ intervals and satisfies
\begin{equation}\label{eq:poly-sublevel-en}
\left|\set{x\in J:\abs{P(x,f(x))}\leq\eps}\right|
\leq C_{d,J}\eps^{1/(N-1)}.
\end{equation}
The constants depend only on $d$, $J$, and the restriction of $f$ to a
neighborhood of $J$.
\end{corollary}

\begin{proof}
Arrange the coefficients $(c_{i0},c_{i1})$ in the same order as the
family used to define $W_d$. Then $P(x,f(x))$ is the corresponding
linear combination of the members of $\mathcal F_d(f)$, and its
coefficient vector has Euclidean norm one. Since $W_d$ is continuous,
nonzero on the compact interval $J$, and hence bounded away from zero,
Lemma~\ref{lem:uniform-sublevel-en} applies and gives the assertion.
\end{proof}

\section{Anisotropic determinant estimates}
\label{sec:determinant}

We compare determinants formed from rational points close to
the graph of $f$. Smoothness bounds these determinants from
above on short source intervals, while clearing the source and
target denominators separately gives a lower bound for every
nonzero determinant. Their comparison will force an auxiliary
relation of degree one in the target variable.

For vectors $v_1,\ldots,v_s\in\R^N$, we equip the exterior power
$\bigwedge^s\R^N$ with its Euclidean norm. Thus
\[
\norm{v_1\wedge\cdots\wedge v_s}^2
=
\det\bigl(\langle v_i,v_j\rangle\bigr)_{1\leq i,j\leq s}.
\]
In particular, when $s=N$,
\[
\norm{v_1\wedge\cdots\wedge v_N}
=
\abs{\det(v_1,\ldots,v_N)}.
\]
We shall also use
\[
\norm{v\wedge\omega}\leq\norm{v}\,\norm{\omega}
\]
for $v\in\R^N$ and $\omega\in\bigwedge^s\R^N$.

\subsection{Exterior decay along a smooth curve}

Several values of a smooth one-dimensional curve taken in an interval
of length $\rho$ become increasingly dependent as $\rho\to0$. The
exponent below reflects the successive Taylor orders
$0,1,\ldots,s-1$.

\begin{lemma}\label{lem:wedge-curve-en}
Let $J\subseteq\R$ be a nondegenerate compact interval, let
$1\leq s\leq N$, and let
\[
\Phi:J\longrightarrow\R^N
\]
be the restriction to $J$ of a $C^s$ map defined on a neighborhood of
$J$. There exists a constant $C_\Phi>0$ such that, whenever
$0<\rho\leq1$ and $t_1,\ldots,t_s$ lie in a common subinterval of $J$
of length at most $\rho$, one has
\begin{equation}\label{eq:wedge-bound-en}
\norm{\Phi(t_1)\wedge\cdots\wedge\Phi(t_s)}
\leq C_\Phi\rho^{s(s-1)/2}.
\end{equation}
\end{lemma}

\begin{proof}
Let
\[
a:=\min_{1\leq i\leq s}t_i,
\qquad
\ell:=\max_{1\leq i\leq s}t_i-a.
\]
If $\ell=0$, the assertion is immediate: for $s\geq2$ the exterior
product vanishes, while for $s=1$ it follows after taking
$C_\Phi\geq\sup_J\norm{\Phi}$. Assume henceforth that
$0<\ell\leq\rho$. Taylor's formula with integral remainder gives
\begin{equation}\label{eq:vector-taylor-en}
\Phi(t_i)
=
\sum_{r=0}^{s-1}(t_i-a)^rv_r+R_i,
\qquad
v_r:=\frac{\Phi^{(r)}(a)}{r!},
\end{equation}
where
\[
R_i
=
\frac{1}{(s-1)!}
\int_a^{t_i}(t_i-u)^{s-1}\Phi^{(s)}(u)\,du.
\]
Since $\Phi^{(s)}$ is bounded on a neighborhood of $J$, there is a
constant $C_0>0$, independent of $a$, $\ell$, and $i$, such that
\begin{equation}\label{eq:remainder-bound-en}
\norm{R_i}\leq C_0\abs{t_i-a}^s\leq C_0\ell^s.
\end{equation}
Also $\abs{t_i-a}^r\leq\ell^r$ for $0\leq r\leq s-1$.

Expand the exterior product of the expressions in
\eqref{eq:vector-taylor-en}. Consider a term containing exactly $k$
remainder vectors. The remaining $s-k$ factors are chosen from
$v_0,\ldots,v_{s-1}$. A term in which some $v_r$ occurs twice vanishes.
Thus, in every nonzero term, the corresponding indices
$r_1,\ldots,r_{s-k}$ are distinct and
\[
r_1+\cdots+r_{s-k}
\geq
0+1+\cdots+(s-k-1)
=
\frac{(s-k)(s-k-1)}{2}.
\]
Such a term is therefore bounded by a constant times
\[
\ell^{ks+(s-k)(s-k-1)/2}.
\]
The identity
\[
ks+\frac{(s-k)(s-k-1)}{2}
=
\frac{s(s-1)}{2}+\frac{k(k+1)}{2}
\]
shows that this is
$O_\Phi(\ell^{s(s-1)/2})$. There are only finitely many terms, and the
vectors $v_0,\ldots,v_{s-1}$ are uniformly bounded as $a$ ranges over
$J$. Hence
\[
\norm{\Phi(t_1)\wedge\cdots\wedge\Phi(t_s)}
\leq C_\Phi\ell^{s(s-1)/2}
\leq C_\Phi\rho^{s(s-1)/2}.
\]
\end{proof}

\subsection{The perturbed determinant}

For $d\geq0$, put $N=2(d+1)$, order the pairs
\[
(i,j)\in\{0,\ldots,d\}\times\{0,1\}
\]
as in Section~\ref{sec:wronskians}, and define
\begin{equation}\label{eq:Psi-d-en}
\Psi_d(x,y):=\bigl(x^iy^j\bigr)_{
\substack{0\leq i\leq d\\j\in\{0,1\}}}\in\R^N,
\qquad
\Phi_d(x):=\Psi_d(x,f(x)).
\end{equation}

\begin{lemma}[Perturbed target-linear determinant]
\label{lem:perturbed-determinant-en}
Let $d\geq0$, put $N=2(d+1)$, and let $J\Subset U$ be a nondegenerate compact
interval. There exists a constant $C_{d,J}^{\mathrm{pert}}>0$, depending
only on $d$, $J$, and the restriction of $f$ to a neighborhood of $J$,
with the following property. Suppose that $0<\rho\leq1$, that
$x_1,\ldots,x_N\in J$ lie in a common subinterval of $J$ of length at
most $\rho$, and that $y_1,\ldots,y_N\in\R$ satisfy
\begin{equation}\label{eq:vertical-error-en}
\abs{y_k-f(x_k)}\leq\delta\leq\rho^N
\qquad(1\leq k\leq N).
\end{equation}
Regarding the vectors $\Psi_d(x_k,y_k)$ as the rows of an
$N\times N$ matrix, one has
\begin{equation}\label{eq:perturbed-simple-en}
\abs{\det\bigl(\Psi_d(x_k,y_k)\bigr)_{1\leq k\leq N}}
\leq
C_{d,J}^{\mathrm{pert}}\rho^{N(N-1)/2}.
\end{equation}
\end{lemma}

\begin{proof}
Set
\[
E_k:=\Psi_d(x_k,y_k)-\Psi_d(x_k,f(x_k)).
\]
The coordinates of $E_k$ corresponding to $j=0$ vanish, whereas those
corresponding to $j=1$ are
\[
x_k^i\bigl(y_k-f(x_k)\bigr),
\qquad0\leq i\leq d.
\]
Since $x_k\in J$, there is a constant $C_0=C_0(d,J)$ such that
\begin{equation}\label{eq:perturbation-vector-en}
\norm{E_k}\leq C_0\delta
\qquad(1\leq k\leq N).
\end{equation}
We have
\[
\Psi_d(x_k,y_k)=\Phi_d(x_k)+E_k.
\]
Expanding the exterior product of these $N$ vectors by multilinearity,
consider a term containing exactly $r$ error vectors. If $r<N$, the
exterior norm inequality, \eqref{eq:perturbation-vector-en}, and
Lemma~\ref{lem:wedge-curve-en}, applied to the remaining $N-r$ curve
vectors, bound this term by
\[
C_r\delta^r\rho^{(N-r)(N-r-1)/2}.
\]
For $r=N$, the same bound holds with the standard convention
$\bigwedge^0\R^N=\R$ and $\norm{1}=1$. Absorbing the finitely many
combinatorial factors, we obtain
\begin{equation}\label{eq:perturbed-general-en}
\abs{\det\bigl(\Psi_d(x_k,y_k)\bigr)_{1\leq k\leq N}}
\leq
C\sum_{r=0}^N
\delta^r\rho^{(N-r)(N-r-1)/2}.
\end{equation}
By \eqref{eq:vertical-error-en}, the summand of index $r$ is at most
\[
\rho^{rN+(N-r)(N-r-1)/2}.
\]
The exact identity
\begin{equation}\label{eq:exponent-identity-en}
rN+\frac{(N-r)(N-r-1)}{2}
=
\frac{N(N-1)}{2}+\frac{r(r+1)}{2}
\end{equation}
shows that, since $0<\rho\leq1$, every summand is at most
$\rho^{N(N-1)/2}$. Summing over $0\leq r\leq N$ proves
\eqref{eq:perturbed-simple-en}.
\end{proof}

\subsection{The arithmetic lower bound}

The preceding estimate is purely analytic. Rationality supplies the
complementary fact that a nonzero target-linear determinant cannot be
too small.

\begin{lemma}\label{lem:det-lower-en}
Fix $d\geq0$ and put $N=2(d+1)$. Let
\[
x_k=\frac{p_k}{q_k},
\qquad
y_k=\frac{a_k}{b_k},
\qquad1\leq k\leq N,
\]
where $p_k,a_k\in\Z$ and $q_k,b_k\in\Z_{>0}$. The fractions need not be
reduced. Define
\[
\Delta
:=
\det\bigl(x_k^iy_k^j\bigr)_{
\substack{1\leq k\leq N\\0\leq i\leq d,\ j\in\{0,1\}}}.
\]
If $\Delta\neq0$, then
\begin{equation}\label{eq:det-lower-product-en}
\abs{\Delta}\geq\prod_{k=1}^Nq_k^{-d}b_k^{-1}.
\end{equation}
In particular, if $q_k<2Q$ and $b_k<2B$ for every $k$, then
\begin{equation}\label{eq:det-lower-dyadic-en}
\abs{\Delta}\geq(2Q)^{-dN}(2B)^{-N}.
\end{equation}
\end{lemma}

\begin{proof}
Multiply the $k$th row of the matrix defining $\Delta$ by $q_k^db_k$.
The entry corresponding to $(i,j)$ becomes
\[
q_k^db_kx_k^iy_k^j
=
\begin{cases}
 p_k^iq_k^{d-i}b_k,&j=0,\\
 p_k^iq_k^{d-i}a_k,&j=1,
\end{cases}
\]
and is therefore an integer. The determinant of the resulting integer
matrix is
\[
\left(\prod_{k=1}^Nq_k^db_k\right)\Delta.
\]
If $\Delta\neq0$, this is a nonzero integer and hence has absolute
value at least one, proving \eqref{eq:det-lower-product-en}. The final assertion follows from $q_k<2Q$ and $b_k<2B$.
\end{proof}

The next section combines the analytic upper bound
\eqref{eq:perturbed-simple-en} with the arithmetic lower bound
\eqref{eq:det-lower-dyadic-en}. The two estimates retain the source
and target heights separately; their balance is achieved by adapting
the source degree to the relative size of these two heights.

\section{Two-height counting on analytic graphs}
\label{sec:scarcity}

We now establish the main counting estimate for the nonrational
real-analytic function $f:U\to\R$ fixed in
Section~\ref{sec:wronskians}. We count rational source points
$p/q$, with $q\asymp Q$, whose images admit exceptionally good
rational approximations, keeping the source and target denominator
scales separate. The numerical parameters below are chosen to
ensure the strict exponent inequalities needed in the proof;
we make no attempt to optimize them.

We fix
\begin{equation}\label{eq:fixed-parameters-en}
\gamma=10,
\qquad
\tau=100,
\qquad
v=\frac{2}{\gamma}=\frac15,
\qquad
\tau-3=97.
\end{equation}
The target degree is one throughout this section. The parameter
$\gamma$ controls the source degree, $\tau$ is the approximation
exponent that will eventually be excluded for the image, and $v$
separates the elementary and determinant ranges of target
denominators.

For every integer $A\geq3$, define
\begin{equation}\label{eq:D-A-en}
D(A)
:=
\max\set{2,\left\lceil\frac{\gamma A}{\tau-3}\right\rceil}
=
\max\set{2,\left\lceil\frac{10A}{97}\right\rceil},
\end{equation}
and, for $Q\geq2$, put
\begin{equation}\label{eq:T-Q-A-en}
T(Q,A):=Q^{A/(\tau-3)}=Q^{A/97}.
\end{equation}
The quantity $A$ measures the required accuracy in the source
variable, whereas $T(Q,A)$ is the largest target height treated at
source scale $Q$. The exponent $\tau-3$ is chosen for the subsequent fusion: for fixed $A$,
$T(Q,A)^{2-\tau}=o(Q^{-A})$ as $Q\to\infty$.
Thus the contribution from the next lower target cutoff can
be made small relative to the next source interval.

Let $I_*\Subset U$ be a nondegenerate compact interval on which
\begin{equation}\label{eq:derivative-bounds-en}
0<m_*\leq\abs{f'(x)}\leq M_*
\qquad(x\in I_*).
\end{equation}
Since $f'$ is continuous and does not vanish on $I_*$, it has constant
sign there. Hence $f|_{I_*}$ is strictly monotone, and the mean-value
theorem gives
\begin{equation}\label{eq:inverse-lipschitz-en}
\abs{x_2-x_1}
\leq
m_*^{-1}\abs{f(x_2)-f(x_1)}
\qquad(x_1,x_2\in I_*).
\end{equation}

For a nondegenerate compact interval $J\subseteq I_*$, integers
$A\geq3$ and $H\geq2$, and an integer $Q\geq2$, define the set of \textit{resonant source centers}
\begin{equation}\label{eq:dangerous-set-en}
\begin{split}
\mathcal D(J;Q,A,H)
:=
\Bigl\{\frac pq\in J:\;&
\gcd(p,q)=1,\quad Q\leq q<2Q,\\
&\exists\,a\in\Z,\ b\in\Z_{>0}
\text{ with }H\leq b<T(Q,A),\\
&\abs{f(p/q)-a/b}
\leq2b^{-\tau}+4M_*Q^{-A}
\Bigr\}.
\end{split}
\end{equation}
The source fractions are reduced, whereas the target fractions need not be.
For $p/q\in\mathcal D(J;Q,A,H)$, any pair $(a,b)$ satisfying the
conditions in \eqref{eq:dangerous-set-en} will be called a \emph{witness}
for $p/q$.
The term $Q^{-A}$ is a stability margin which will allow avoidance at a
rational center to persist on the next source interval.

\begin{proposition}[Two-height counting on analytic graphs]
\label{prop:scarcity-en}
There exists a constant $C_{\mathrm{low}}>0$, depending only on
$f$, $I_*$, $m_*$, $M_*$ and the fixed parameters in
\eqref{eq:fixed-parameters-en}, with the following property. Let
$J\subseteq I_*$ be a nondegenerate compact interval and let $A\geq3$
be an integer. Assume that
\begin{equation}\label{eq:wronskians-on-J-en}
W_d(x)\neq0
\qquad(x\in J,\ 2\leq d\leq D(A)).
\end{equation}
Then there exist constants
\[
C(J,A)>0
\qquad\text{and}\qquad
Q_0(J,A)\in\Z_{>0}
\]
such that, for every integer $H\geq2$ and every integer
$Q\geq Q_0(J,A)$,
\begin{equation}\label{eq:scarcity-bound-en}
\#\mathcal D(J;Q,A,H)
\leq
C_{\mathrm{low}}Q^2H^{-98}
+
C(J,A)Q^{17/10}.
\end{equation}
In particular, the constants $C(J,A)$ and $Q_0(J,A)$ are uniform in
$H$.
\end{proposition}

\begin{proof}
Fix $J$ and $A$ as in the statement. Every lower bound imposed on $Q$
below will depend at most on $J$ and $A$, and will be independent of
$H$ and of the target block.

If $T(Q,A)\leq H$, then
$\mathcal D(J;Q,A,H)=\varnothing$. Assume henceforth that
$T(Q,A)>H$. For every integer $k\geq0$ such that
\[
B_k:=2^kH<T(Q,A),
\]
put
\begin{equation}\label{eq:dyadic-blocks-en}
\mathcal B_k
:=
[B_k,\min\{2B_k,T(Q,A)\})\cap\Z_{>0}.
\end{equation}
These sets form a disjoint partition of
$[H,T(Q,A))\cap\Z_{>0}$. Let $\mathcal D_k$ be the set of source centers
admitting at least one
witness $(a,b)$ with
$b\in\mathcal B_k$. Then
\begin{equation}\label{eq:D-covered-by-blocks-en}
\mathcal D(J;Q,A,H)\subseteq\bigcup_k\mathcal D_k,
\qquad
\#\mathcal D(J;Q,A,H)\leq\sum_k\#\mathcal D_k.
\end{equation}
A center may occur in more than one $\mathcal D_k$, which only produces
an admissible overcount.

A block with left endpoint $B=B_k$ is called \emph{small} if
$B<Q^v$ and \emph{large} if $B\geq Q^v$. A block beginning below
$Q^v$ and crossing that threshold is assigned to the small range; the
estimates there use only $B\leq b<2B$.

\subsubsection*{Small target heights}

Set
\begin{equation}\label{eq:F-star-C-star-en}
F_*:=\max_{x\in I_*}\abs{f(x)},
\qquad
C_*:=F_*+2^{1-\tau}+\frac{M_*}{2}.
\end{equation}
Suppose that $x\in I_*$, $b\geq2$, $A\geq3$, $Q\geq2$, and
\[
\abs{f(x)-a/b}\leq2b^{-\tau}+4M_*Q^{-A}.
\]
Since $Q^{-A}\leq2^{-3}$,
\[
\frac{|a|}{b}
\leq
F_*+2b^{-\tau}+4M_*Q^{-A}
\leq C_*.
\]
Thus $|a|\leq C_*b$. For $B\leq b<2B$, there are therefore
$O(B^2)$ pairs $(a,b)$ that can occur as witnesses, with an implied constant depending only on $C_*$.

For each such pair, let
\[
E_{a,b}
:=
\set{x\in J:
\abs{f(x)-a/b}\leq2b^{-\tau}+4M_*Q^{-A}}.
\]
Because $f$ is strictly monotone on $I_*$, the set $E_{a,b}$ is empty,
a singleton, or an interval. If $x_1<x_2$ belong to $E_{a,b}$, then
\eqref{eq:inverse-lipschitz-en} gives
\[
m_*(x_2-x_1)
\leq
\abs{f(x_2)-f(x_1)}
\leq
2\bigl(2b^{-\tau}+4M_*Q^{-A}\bigr).
\]
Consequently,
\begin{equation}\label{eq:Eab-length-en}
|E_{a,b}|
\leq
\frac{2}{m_*}\bigl(2b^{-\tau}+4M_*Q^{-A}\bigr).
\end{equation}

Let $E_B$ be the union of the sets $E_{a,b}$ over one small target
block. It is a union of $O(B^2)$ intervals, and
\begin{equation}\label{eq:low-target-length-en}
|E_B|\ll B^{2-\tau}+B^2Q^{-A}.
\end{equation}
Every center in the corresponding $\mathcal D_k$ belongs to $E_B$.
Applying Lemma~\ref{lem:rational-union} gives
\begin{align}
\#\mathcal D_k
&\ll Q^2|E_B|+B^2\notag\\
&\ll
Q^2B^{2-\tau}
+Q^{2-A}B^2
+B^2.
\label{eq:low-block-count-en}
\end{align}

We sum over the dyadic left endpoints $B=2^kH<Q^v$. Since
$2-\tau<0$,
\[
\sum_{2^kH<Q^v}(2^kH)^{2-\tau}
\ll H^{2-\tau},
\]
whereas
\[
\sum_{2^kH<Q^v}(2^kH)^2\ll Q^{2v}.
\]
It follows that
\begin{equation}\label{eq:low-total-en}
\sum_{B_k<Q^v}\#\mathcal D_k
\leq
C_{\mathrm{low}}Q^2H^{2-\tau}
+O\bigl(Q^{2-A+2v}+Q^{2v}\bigr).
\end{equation}
The coefficient $C_{\mathrm{low}}$ depends only on
$f$, $I_*$, $m_*$, $M_*$ and the fixed parameters in
\eqref{eq:fixed-parameters-en}. Since $A\geq3$ and $v=1/5$,
\[
2-A+2v\leq-\frac35,
\qquad
2v=\frac25.
\]
Thus
\begin{equation}\label{eq:low-final-en}
\sum_{B_k<Q^v}\#\mathcal D_k
\leq
C_{\mathrm{low}}Q^2H^{2-\tau}+O(Q^{2/5}).
\end{equation}

\subsubsection*{Large target heights}

If there are no blocks with $B_k\geq Q^v$, there is nothing to prove. This is automatically the case when $A\leq 19$, since then
\[
T(Q,A)=Q^{A/97}<Q^{1/5}=Q^v.
\]
Hence, in the large-target range one necessarily has $A\geq 20$. Fix one such block, write $B=B_k$, and define
\begin{equation}\label{eq:u-definition-en}
u:=\frac{\log B}{\log Q},
\qquad B=Q^u.
\end{equation}
Since $Q^v\leq B<T(Q,A)$,
\begin{equation}\label{eq:u-range-en}
\frac15=v\leq u<\frac{A}{97}.
\end{equation}
We shall show that each large target block contributes 
$O_{J,A}(Q^{33/20})$ resonant centers: we partition $J$ into
$O_J(Q^\kappa)$ short source cells and prove that each cell contains
only $O_{J,A}(1)$ resonant centers.
Set
\begin{equation}\label{eq:d-N-kappa-en}
d:=\lceil\gamma u\rceil=\lceil10u\rceil,
\qquad
N:=2(d+1),
\qquad
\kappa:=\frac{3(d+u)}{N-1},
\end{equation}
and put
\begin{equation}\label{eq:rho-definition-en}
\rho:=Q^{-\kappa}.
\end{equation}
The degree $d\asymp u$ adapts the source degree to the relative target
height $B=Q^u$, while the choice of $\kappa$ gives
\[
\rho^{N(N-1)/2}=Q^{-\frac32N(d+u)},
\]
to be compared below with the arithmetic scale $Q^{-N(d+u)}$.
Because $10u\geq2$, we have $d\geq2$. Moreover,
$10u<10A/97$, and hence $d\leq D(A)$. Thus
\eqref{eq:wronskians-on-J-en} applies to the chosen degree.

Partition $J$ into half-open cells of length at most $\rho$, taking the
final cell to be closed at its right endpoint. Every point of $J$ then
belongs to exactly one cell, and the number of cells is at most
\begin{equation}\label{eq:number-cells-en}
1+|J|\rho^{-1}
=1+|J|Q^\kappa
\leq(1+|J|)Q^\kappa.
\end{equation}

Fix a cell $C$. For every center $x=p/q\in\mathcal D_k\cap C$, choose once and for all
a witness $(a,b)$ with $b\in\mathcal B_k$, and put $y=a/b$. Define
\begin{equation}\label{eq:s-definition-en}
s:=\min\{u\tau,A\}.
\end{equation}
Since $b\geq B=Q^u$,
\begin{align}
\abs{y-f(x)}
&\leq2B^{-\tau}+4M_*Q^{-A}\notag\\
&\leq K_0Q^{-s},
\qquad K_0:=2+4M_*.
\label{eq:vertical-Q-s-en}
\end{align}

We record uniform exponent bounds. From $d=\lceil10u\rceil$ and
$u\geq1/5$,
\begin{equation}\label{eq:basic-d-bounds-en}
10u\leq d\leq10u+1\leq15u,
\qquad
d+u\leq16u.
\end{equation}
Since $u\leq d/10$ and $N-1=2d+1$,
\begin{equation}\label{eq:kappa-bound-en}
\kappa
=\frac{3(d+u)}{2d+1}
\leq\frac{3(11d/10)}{2d}
=\frac{33}{20}.
\end{equation}
The upper bound in \eqref{eq:u-range-en} gives $97u<A$, while
$\tau u=100u$. Therefore
\begin{equation}\label{eq:s-lower-en}
s\geq97u.
\end{equation}
Also, using $N/(N-1)\leq2$ and
\eqref{eq:basic-d-bounds-en},
\begin{equation}\label{eq:kappaN-upper-en}
\kappa N
=3(d+u)\frac{N}{N-1}
\leq6(d+u)
\leq96u.
\end{equation}
Consequently,
\begin{equation}\label{eq:s-kappa-gap-en}
s-\kappa N\geq u\geq\frac{1}{5}.
\end{equation}
Finally,
\begin{equation}\label{eq:N-upper-en}
N-1=2d+1\leq30u+1\leq35u.
\end{equation}
Thus the three numerical inequalities driving the large-target
argument are
\[
\kappa\leq\frac{33}{20},\qquad
s-\kappa N\geq\frac15,\qquad
\frac{s}{N-1}\geq\frac{97}{35}>2.
\]
They control, respectively, the number of source cells, the
perturbation required in Lemma~\ref{lem:perturbed-determinant-en},
and the Farey count within each cell.

\emph{Determinant comparison.} For each integer $d$ with $2\leq d\leq D(A)$, let
$C_{d,J}^{\mathrm{pert}}$ be the constant in
Lemma~\ref{lem:perturbed-determinant-en}, and set
\begin{equation}\label{eq:uniform-det-constant-en}
K_{\det}(J,A)
:=
\max_{2\leq d\leq D(A)}C_{d,J}^{\mathrm{pert}}<\infty.
\end{equation}
Since $D(A)$ is fixed once $A$ is fixed, this maximum is taken over
a finite set of constants determined before the choice of $Q$ and of
the target block. In particular, $K_{\det}(J,A)$ is independent of
$Q$, $u$, $B$, and $H$.

By \eqref{eq:vertical-Q-s-en} and
\eqref{eq:s-kappa-gap-en},
\[
K_0Q^{-s}\leq K_0Q^{-\kappa N-1/5}.
\]
Hence $Q\geq K_0^5$ implies
\[
K_0Q^{-s}\leq Q^{-\kappa N}=\rho^N.
\]
This restriction is simultaneous for all large target blocks and is
independent of $u$, $B$, and $H$.

Suppose first that $C$ contains at least $N$ resonant centers. Choose
$N$ distinct centers and their fixed witnessing pairs, say
$(x_r,y_r)$, $1\leq r\leq N$, and define
\[
\Delta
:=
\det\bigl(x_r^iy_r^j\bigr)_{
\substack{1\leq r\leq N\\0\leq i\leq d,\ j\in\{0,1\}}}.
\]
All the $x_r$ lie in the closure of a cell of length at most $\rho$.
Lemma~\ref{lem:perturbed-determinant-en} and the definition of
$\kappa$ give
\begin{equation}\label{eq:det-upper-high-en}
|\Delta|
\leq
K_{\det}(J,A)\rho^{N(N-1)/2}
=
K_{\det}(J,A)Q^{-\frac32N(d+u)}.
\end{equation}
If $\Delta\neq0$, then $q_r<2Q$ and $b_r<2B$, so
Lemma~\ref{lem:det-lower-en} gives
\begin{equation}\label{eq:det-lower-high-en}
|\Delta|
\geq
(2Q)^{-dN}(2B)^{-N}
=
2^{-N(d+1)}Q^{-N(d+u)}.
\end{equation}
The two bounds are incompatible whenever
\begin{equation}\label{eq:det-incompatibility-en}
Q^{N(d+u)/2}
>
K_{\det}(J,A)2^{N(d+1)}.
\end{equation}
This threshold is uniform in the target block. Indeed,
$d\geq2$, $u\geq1/5$, and $N\geq6$, so
\[
\frac{N(d+u)}{2}
\geq\frac{6(2+1/5)}{2}
=\frac{33}{5}>6.
\]
Since $N(d+1)=2(d+1)^2$, it is enough to impose
\[
Q^6
>
K_{\det}(J,A)
\max_{2\leq d\leq D(A)}2^{2(d+1)^2}.
\]
This condition depends only on $J$ and $A$. Thus, for all sufficiently
large $Q$, the determinant associated with every choice of $N$
distinct resonant centers in every cell is zero.

Let $M_C$ be the number of resonant centers in $C$ for the present
target block. If $M_C<N$, then
\[
M_C<N=2(d+1)\leq2(D(A)+1),
\]
so $M_C=O_A(1)$. Assume that $M_C\ge N$ and enumerate the centers together with their
chosen witnesses as
\[
(x_r,y_r),\qquad 1\le r\le M_C,
\]
where $x_r=p_r/q_r$ and $y_r=a_r/b_r$. Consider the $M_C\times N$ evaluation matrix
\begin{equation}\label{eq:evaluation-matrix-en}
\mathcal M_C
:=
\bigl(x_r^iy_r^j\bigr)_{
\substack{1\leq r\leq M_C\\0\leq i\leq d,\ j\in\{0,1\}}}.
\end{equation}
Since the preceding choice of $N$ distinct centers was arbitrary, every
$N\times N$ minor of $\mathcal M_C$ vanishes. Hence
$\rank\mathcal M_C\leq N-1$, so the right kernel of $\mathcal M_C$ in
$\R^N$ is nontrivial. Choose a unit vector $(c_{i0},c_{i1})$ in this
right kernel and define
\begin{equation}\label{eq:auxiliary-polynomial-en}
P_C(X,Y)
:=
A_C(X)+B_C(X)Y
:=
\sum_{i=0}^dc_{i0}X^i
+
\sum_{i=0}^dc_{i1}X^iY,
\qquad
\norm{P_C}_2=1.
\end{equation}
Then
\begin{equation}\label{eq:P-vanishes-pairs-en}
P_C(x_r,y_r)=0
\qquad(1\leq r\leq M_C).
\end{equation}

\emph{Counting in one cell.} Put
\[
X_*:=\max\set{1,\max_{x\in J}|x|}
\]
and
\begin{equation}\label{eq:L-A-en}
L_{J,A}
:=
\max_{2\leq d\leq D(A)}
\left(\sum_{i=0}^dX_*^{2i}\right)^{1/2}<\infty.
\end{equation}
The coefficient normalization in
\eqref{eq:auxiliary-polynomial-en} and Cauchy--Schwarz give
$|B_C(x)|\leq L_{J,A}$ for $x\in J$. Since $P_C$ is linear in $Y$,
\eqref{eq:P-vanishes-pairs-en} and
\eqref{eq:vertical-Q-s-en} imply
\begin{equation}\label{eq:P-small-graph-en}
\abs{P_C(x_r,f(x_r))}
=
|B_C(x_r)|\abs{f(x_r)-y_r}
\leq
L_{J,A}K_0Q^{-s}.
\end{equation}

For $2\leq d\leq D(A)$, let
$C_{d,J}$, $R_{d,J}$, and $\eps_{d,J}$ be supplied by
Corollary~\ref{cor:polynomial-sublevel-en}, and set
\begin{equation}\label{eq:uniform-sublevel-data-en}
\begin{aligned}
C_{J,A}^{\mathrm{sub}}
&:=\max_{2\leq d\leq D(A)}C_{d,J},\\
R_{J,A}^{\mathrm{sub}}
&:=\max_{2\leq d\leq D(A)}R_{d,J},\\
\eps_{J,A}^{\mathrm{sub}}
&:=\min_{2\leq d\leq D(A)}\eps_{d,J}>0.
\end{aligned}
\end{equation}
These extrema are finite. Since $s\geq97/5$,
\[
L_{J,A}K_0Q^{-s}\leq L_{J,A}K_0Q^{-97/5},
\]
so, after increasing the lower threshold for $Q$ depending only on
$J$ and $A$, we have
\begin{equation}\label{eq:sublevel-smallness-en}
L_{J,A}K_0Q^{-s}<\eps_{J,A}^{\mathrm{sub}}.
\end{equation}
This threshold is again uniform over all large target blocks.
Define
\[
S_C
:=
\set{x\in J:
\abs{P_C(x,f(x))}\leq L_{J,A}K_0Q^{-s}}
\cap\overline C.
\]
Intersecting the sublevel set from
Corollary~\ref{cor:polynomial-sublevel-en} with the interval
$\overline C$ does not increase its number of connected components.
Every resonant center assigned to $C$ belongs to $S_C$. By
Corollary~\ref{cor:polynomial-sublevel-en},
\begin{equation}\label{eq:sublevel-length-high-en}
|S_C|\ll_{J,A}Q^{-s/(N-1)},
\end{equation}
and $S_C$ is a union of at most
$R_{J,A}^{\mathrm{sub}}$ intervals. From
\eqref{eq:s-lower-en} and \eqref{eq:N-upper-en},
\begin{equation}\label{eq:s-over-N-en}
\frac{s}{N-1}\geq\frac{97}{35}>2.
\end{equation}
The resonant centers are distinct reduced fractions with denominators
in $[Q,2Q)$. By \eqref{eq:sublevel-length-high-en} and
\eqref{eq:s-over-N-en},
\[
Q^2|S_C|
\ll_{J,A}
Q^{2-s/(N-1)}
\leq Q^{2-97/35}
=Q^{-27/35}.
\]
Lemma~\ref{lem:rational-union} therefore gives
\begin{align}
M_C
&\leq4Q^2|S_C|+R_{J,A}^{\mathrm{sub}}\notag\\
&\ll_{J,A}Q^{-27/35}+1
\ll_{J,A}1.
\label{eq:centers-per-cell-en}
\end{align}
Together with the case $M_C<N$, this bound holds uniformly for every
cell and every large target block.

\emph{Summation over cells and target blocks.} By \eqref{eq:number-cells-en}, there are $O_J(Q^\kappa)$ source
cells, and \eqref{eq:centers-per-cell-en} gives $O_{J,A}(1)$ resonant
centers in each. Hence, by \eqref{eq:kappa-bound-en}, one large target
block contributes
\begin{equation}\label{eq:high-one-block-en}
O_{J,A}(Q^\kappa)=O_{J,A}(Q^{33/20}).
\end{equation}
The total number of dyadic target blocks is at most
\[
1+\log_2\frac{T(Q,A)}{H}
\leq
1+\frac{A}{97}\log_2Q
=O_A(\log Q).
\]
Since $\log Q\ll Q^{1/20}$ for $Q\geq2$, summing
\eqref{eq:high-one-block-en} gives
\begin{equation}\label{eq:high-all-blocks-en}
\sum_{B_k\geq Q^v}\#\mathcal D_k
=O_{J,A}(Q^{33/20}\log Q)
=O_{J,A}(Q^{17/10}).
\end{equation}

Combining \eqref{eq:D-covered-by-blocks-en},
\eqref{eq:low-final-en}, and \eqref{eq:high-all-blocks-en}, we obtain
\eqref{eq:scarcity-bound-en}. We finally record the uniformity. For
fixed $J$ and $A$, only the finite range $2\leq d\leq D(A)$ occurs,
and the bounds $u\geq1/5$ and $s-\kappa N\geq1/5$ make all lower
restrictions on $Q$ in the large-height argument uniform in the target
block. Thus none depends on $H$, $u$, $B$, the cell, or the witnessing
pairs. Taking their maximum defines $Q_0(J,A)$, and enlarging the final
implied constant defines $C(J,A)$.
\end{proof}

%\begin{remark}\label{rem:anisotropy}
%The target degree one is essential to the conclusion. The determinant
%produces a relation $A_C(X)+B_C(X)Y=0$, while the Wronskian condition
%asserts that the analytic graph is not itself rational. The sharp
%counting input is Lemma~\ref{lem:rational-union}: Farey separation
%removes the factor $Q$ per connected component that would otherwise
%force a higher target degree. The source degree
%$d=\lceil10\log B/\log Q\rceil$ balances the two independent heights.
%For fixed $A$, only the finite range $2\leq d\leq D(A)$ occurs, so all
%Wronskian, determinant, and sublevel constants are fixed before
%$Q\to\infty$.
%\end{remark}

\section{From two-height counting to rigidity}
\label{sec:fusion}

\subsection{A safe center}

The counting estimate becomes effective once its quadratic main term
is smaller than the supply of reduced source fractions. For the
deduction below, the precise exponent $17/10$ is not essential: what
matters is that the second term in Proposition~5.1 is $o(Q^2)$ as
$Q\to\infty$, uniformly in the lower target cutoff $H$, while
$H^{2-\tau}\to0$ as $H\to\infty$. These two features produce safe
rational centers at every sufficiently large source scale.

\begin{lemma}[A safe center in one source block]
\label{lem:safe-center-en}
Assume the hypotheses and notation of
Proposition~\ref{prop:scarcity-en}. Let $H\geq2$ be an integer and
suppose in addition that
\begin{equation}\label{eq:tail-smallness-en}
C_{\mathrm{low}}H^{2-\tau}
\leq
\frac{c_{\mathrm F}}4|J|.
\end{equation}
Then, for every sufficiently large integer $Q$, there exists a reduced
fraction
\[
r=\frac pq\in J,
\qquad Q\leq q<2Q,
\]
such that
\begin{equation}\label{eq:safe-center-en}
\abs{f(r)-a/b}
>
2b^{-\tau}+4M_*Q^{-A}
\end{equation}
for every $a\in\Z$ and every integer $b$ satisfying
\[
H\leq b<T(Q,A).
\]
\end{lemma}

\begin{proof}
By Lemma~\ref{lem:farey-lower-en}, for every sufficiently large $Q$,
the interval $J$ contains at least
\[
c_{\mathrm F}|J|Q^2
\]
reduced fractions with denominator in $[Q,2Q)$. On the other hand,
Proposition~\ref{prop:scarcity-en} and
\eqref{eq:tail-smallness-en} give
\[
\#\mathcal D(J;Q,A,H)
\leq
\frac{c_{\mathrm F}}4|J|Q^2+C(J,A)Q^{17/10}.
\]
Since $17/10<2$ and $J,A$ are fixed, there exists
$Q_1(J,A)$ such that, for every $Q\geq Q_1(J,A)$,
\[
C(J,A)Q^{17/10}
\leq
\frac{c_{\mathrm F}}4|J|Q^2.
\]
Thus $\mathcal D(J;Q,A,H)$ does not exhaust the reduced fractions
$p/q\in J$ with $Q\leq q<2Q$. Hence any such fraction outside
$\mathcal D(J;Q,A,H)$ satisfies \eqref{eq:safe-center-en} by definition.
\end{proof}

\subsection{Fusion and proof of Theorem~\ref{thm:intro-local}}

Fix a nonempty open interval $V\subseteq U$. Since $f$ is not rational, it is nonconstant, so $f'\not\equiv0$. As $f'$ is real-analytic, its zeros are isolated. We may therefore choose a
nondegenerate compact interval
\begin{equation}\label{eq:I-star-choice-en}
I_*\Subset V
\end{equation}
on which $f'$ does not vanish. By continuity and compactness, there are
constants $m_*,M_*>0$ such that
\[
0<m_*\leq\abs{f'(x)}\leq M_*
\qquad(x\in I_*),
\]
as required in \eqref{eq:derivative-bounds-en}.

For $n\geq0$, put
\begin{equation}\label{eq:A-n-en}
A_n:=n+3
\end{equation}
and define
\begin{equation}\label{eq:Z-n-en}
\mathcal Z_n
:=
\bigcup_{2\leq d\leq D(A_n)}
\set{x\in I_*:W_d(x)=0}.
\end{equation}
Each $\mathcal Z_n$ is finite by
Corollary~\ref{cor:wronskian-finite-en}. Since $D(A)$ is nondecreasing,
\begin{equation}\label{eq:Z-n-nested-en}
\mathcal Z_n\subseteq\mathcal Z_{n+1}.
\end{equation}
Choose a nondegenerate compact interval
\begin{equation}\label{eq:I0-en}
I_0\subseteq\Int(I_*)\setminus\mathcal Z_0.
\end{equation}

For a nondegenerate compact interval $I=[\alpha,\beta]$, let
\begin{equation}\label{eq:middle-third-en}
J(I)
:=
\left[
\alpha+\frac{\beta-\alpha}{3},
\beta-\frac{\beta-\alpha}{3}
\right]
\end{equation}
be its middle third. Whenever $I_n$ is defined, write
$J_n:=J(I_n)$. Then
\begin{equation}\label{eq:middle-third-length-en}
|J_n|=\frac13|I_n|,
\qquad
\dist(J_n,\R\setminus I_n)=\frac13|I_n|.
\end{equation}

Choose an integer $T_0\geq2$ such that
\begin{equation}\label{eq:initial-tail-en}
C_{\mathrm{low}}T_0^{2-\tau}
\leq
\frac{c_{\mathrm F}}4|J_0|.
\end{equation}
This is possible because $2-\tau=-98<0$.

We construct nondegenerate compact intervals $I_n$, strictly
increasing positive integers $T_n$, and reduced fractions
$r_n=p_n/q_n$ satisfying the following properties whenever the objects
involved are defined:
\begin{enumerate}[label=\textup{(F\arabic*)},leftmargin=3em]
\item $I_{n+1}\subset\Int(I_n)$;
\item $I_n\cap\mathcal Z_n=\varnothing$;
\item $T_{n+1}>T_n$;
\item $q_n>2q_{n-1}$ for $n\geq1$;
\item for every $x\in I_{n+1}$,
\begin{equation}\label{eq:source-invariant-en}
0<\abs{x-r_n}<q_n^{-A_n};
\end{equation}
\item for every $x\in I_{n+1}$, every $a\in\Z$, and every integer $b$
with $T_n\leq b<T_{n+1}$,
\begin{equation}\label{eq:target-invariant-en}
\abs{f(x)-a/b}>b^{-\tau};
\end{equation}
\item
\begin{equation}\label{eq:tail-invariant-en}
C_{\mathrm{low}}T_n^{2-\tau}
\leq
\frac{c_{\mathrm F}}4|J_n|.
\end{equation}
\end{enumerate}
Conditions (F5) and (F6) impose source approximation and target
avoidance. Conditions (F2) and (F7) retain the hypotheses needed
to choose the next safe center.

For clarity, we write \textup{(F1)}$_n$, \textup{(F3)}$_n$,
\textup{(F5)}$_n$, and \textup{(F6)}$_n$ for the corresponding
properties linking stage $n$ to stage $n+1$, while
\textup{(F2)}$_n$ and \textup{(F7)}$_n$ refer to the properties at
stage $n$ itself; similarly, \textup{(F4)}$_n$ is understood for
$n\geq1$.

The initial choices $I_0$ and $T_0$ satisfy
\textup{(F2)}$_0$ and \textup{(F7)}$_0$ by
\eqref{eq:I0-en} and \eqref{eq:initial-tail-en}. Suppose that $I_0,\ldots,I_n$, $T_0,\ldots,T_n$, and, when
$n\geq1$, $r_0,\ldots,r_{n-1}$ have been constructed and satisfy
all applicable induction hypotheses. We choose
$Q_n,r_n,T_{n+1},I_{n+1}$ in this order and verify the required
conditions at the next stage.

Since $J_n\subseteq I_n$ and
\textup{(F2)}$_n$ gives $I_n\cap\mathcal Z_n=\varnothing$,
\[
W_d(x)\neq0
\qquad(x\in J_n,\ 2\leq d\leq D(A_n)).
\]
Thus Proposition~\ref{prop:scarcity-en} applies on $J_n$ with
$A=A_n$, and \textup{(F7)}$_n$ is precisely the additional hypothesis
\eqref{eq:tail-smallness-en} of Lemma~\ref{lem:safe-center-en}, with
$H=T_n$.

Before choosing the source height, define
\begin{equation}\label{eq:Lambda-next-en}
\Lambda_{n+1}:=\#\mathcal Z_{n+1}+2.
\end{equation}
Notice that $\mathcal Z_{n+1}$, and hence $\Lambda_{n+1}$, is already
fixed at this point, independently of all objects to be chosen at
stage $n$.

Choose an integer $Q_n$ sufficiently large that the conclusion of
Lemma~\ref{lem:safe-center-en} holds and, simultaneously,
\begin{enumerate}[label=\textup{(Q\arabic*)},leftmargin=3em]
\item $Q_n>2q_{n-1}$ if $n\geq1$;
\item
\[
(2Q_n)^{-A_n}<\frac{|I_n|}{3};
\]
\item
\[
Q_n^{A_n/97}\geq2T_n+2;
\]
\item
\begin{equation}\label{eq:Q4-en}
3C_{\mathrm{low}}\Lambda_{n+1}
2^{A_n+\tau-2}Q_n^{-A_n/97}
\leq
\frac{c_{\mathrm F}}4.
\end{equation}
\end{enumerate}
Such a choice is possible. Indeed, at this stage $J_n$, $A_n$,
$T_n$, $\Lambda_{n+1}$, and, when $n\geq1$, $q_{n-1}$ are fixed, as
are all constants and thresholds in
Lemma~\ref{lem:safe-center-en}. Conditions
\textup{(Q1)}--\textup{(Q3)} impose only finitely many lower bounds on
$Q_n$, while the left-hand side of \eqref{eq:Q4-en} tends to zero as
$Q_n\to\infty$.

Choose a reduced safe fraction
\begin{equation}\label{eq:r-n-choice-en}
r_n=\frac{p_n}{q_n}\in J_n,
\qquad Q_n\leq q_n<2Q_n.
\end{equation}
Set
\begin{equation}\label{eq:R-T-next-en}
R_n:=(2Q_n)^{-A_n},
\qquad
T_{n+1}:=\left\lfloor Q_n^{A_n/97}\right\rfloor,
\end{equation}
and define
\begin{equation}\label{eq:K-next-en}
K_{n+1}:=(r_n-R_n,r_n+R_n).
\end{equation}
Condition \textup{(Q3)} gives $T_{n+1}\geq2T_n+2>T_n$, proving
\textup{(F3)}$_n$. Since $r_n\in J_n$ and
$\dist(J_n,\R\setminus I_n)=|I_n|/3$, condition \textup{(Q2)} gives
\[
K_{n+1}\subset\Int(I_n).
\]

We transfer the safety of $r_n$ to $K_{n+1}$. If
$T_n\leq b<T_{n+1}$, then
\[
b<T_{n+1}\leq Q_n^{A_n/97}=T(Q_n,A_n).
\]
Therefore the safety of $r_n$ gives, for every $a\in\Z$,
\begin{equation}\label{eq:safe-at-center-en}
\abs{f(r_n)-a/b}
>
2b^{-\tau}+4M_*Q_n^{-A_n}.
\end{equation}
For $x\in K_{n+1}$, the mean-value theorem and
\eqref{eq:derivative-bounds-en} yield
\[
\abs{f(x)-f(r_n)}
\leq M_*\abs{x-r_n}
<M_*R_n
\leq M_*Q_n^{-A_n}.
\]
Consequently,
\begin{equation}\label{eq:avoidance-on-K-en}
\abs{f(x)-a/b}
>
2b^{-\tau}+3M_*Q_n^{-A_n}
>b^{-\tau}.
\end{equation}

It remains to choose $I_{n+1}$ inside $K_{n+1}$. For every $k<n$,
property \textup{(F5)}$_k$ gives
$r_k\notin I_{k+1}$, and repeated use of \textup{(F1)} gives
$I_n\subseteq I_{k+1}$. Hence no previously chosen center $r_k$,
$k<n$, lies in $K_{n+1}$. It is therefore enough to remove
\begin{equation}\label{eq:forbidden-next-en}
F_{n+1}:=\mathcal Z_{n+1}\cup\{r_n\}.
\end{equation}
The set $F_{n+1}\cap K_{n+1}$ has at most
$\#\mathcal Z_{n+1}+1=\Lambda_{n+1}-1$ points. Hence
$K_{n+1}\setminus F_{n+1}$ has at most $\Lambda_{n+1}$ connected
components. Their total length is $2R_n$, so one component, say
$(u_{n+1},v_{n+1})$, has length at least
$2R_n/\Lambda_{n+1}$. Put
\[
s_n:=\frac{R_n}{\Lambda_{n+1}}
\]
and choose
\begin{equation}\label{eq:I-next-choice-en}
I_{n+1}
:=
\left[
 u_{n+1}+\frac{s_n}{2},
 u_{n+1}+\frac{3s_n}{2}
\right].
\end{equation}
Then $I_{n+1}\Subset(u_{n+1},v_{n+1})$ and
\begin{equation}\label{eq:I-next-length-en}
|I_{n+1}|
=s_n
=\frac{R_n}{\Lambda_{n+1}}
=\frac{(2Q_n)^{-A_n}}{\Lambda_{n+1}}.
\end{equation}
Thus \textup{(F1)}$_n$ and \textup{(F2)}$_{n+1}$ hold.
Furthermore, $r_n\notin I_{n+1}$ and
$I_{n+1}\subset K_{n+1}$, so for every $x\in I_{n+1}$,
\[
0<\abs{x-r_n}<R_n\leq q_n^{-A_n},
\]
where the last inequality follows from $q_n<2Q_n$. This proves
\textup{(F5)}$_n$. Property \textup{(F6)}$_n$ follows from
\eqref{eq:avoidance-on-K-en}. If $n\geq1$, then
$q_n\geq Q_n>2q_{n-1}$ by \textup{(Q1)}, proving
\textup{(F4)}$_n$.

It remains to verify \textup{(F7)}$_{n+1}$. From
\eqref{eq:I-next-length-en} and
\eqref{eq:middle-third-length-en},
\begin{equation}\label{eq:J-next-lower-en}
|J_{n+1}|
=\frac{(2Q_n)^{-A_n}}{3\Lambda_{n+1}}.
\end{equation}
Condition \textup{(Q3)} implies $Q_n^{A_n/97}\geq2$. Since
$\lfloor X\rfloor\geq X/2$ for $X\geq2$,
\[
T_{n+1}\geq\frac12Q_n^{A_n/97}.
\]
Because $2-\tau<0$,
\begin{equation}\label{eq:T-tail-upper-en}
T_{n+1}^{2-\tau}
\leq
2^{\tau-2}Q_n^{-A_n(\tau-2)/97}.
\end{equation}
Combining \eqref{eq:J-next-lower-en} and \eqref{eq:T-tail-upper-en} gives
\begin{align}
\frac{C_{\mathrm{low}}T_{n+1}^{2-\tau}}{|J_{n+1}|}
&\leq
3C_{\mathrm{low}}\Lambda_{n+1}
2^{A_n+\tau-2}
Q_n^{A_n-A_n(\tau-2)/97}\notag\\
&=
3C_{\mathrm{low}}\Lambda_{n+1}
2^{A_n+\tau-2}Q_n^{-A_n/97}\notag\\
&\leq\frac{c_{\mathrm F}}4,
\label{eq:tail-ratio-next-en}
\end{align}
where we used $97=\tau-3$ and \textup{(Q4)}. The negative exponent $-A_n/97$ ensures that the contribution
of the next lower target cutoff becomes small relative to the
length of the next source interval. This is
\textup{(F7)}$_{n+1}$. We have therefore completed the inductive step,
and hence the construction.

By \textup{(F1)}, the intervals $I_n$ are nested, nonempty, and
compact. Moreover, for $n\geq1$,
\[
Q_n>2q_{n-1}\geq2Q_{n-1},
\]
so $Q_n\to\infty$. Since $I_{n+1}\subset K_{n+1}$,
\[
\diam(I_{n+1})\leq2R_n=2(2Q_n)^{-A_n}\longrightarrow0.
\]
The nested-interval theorem therefore gives a unique point
\begin{equation}\label{eq:xi-intersection-en}
\set{\xi}=\bigcap_{n=0}^{\infty}I_n.
\end{equation}
Property \textup{(F5)} yields
\begin{equation}\label{eq:xi-rn-en}
0<\abs{\xi-r_n}<q_n^{-A_n}
\qquad(n\geq0).
\end{equation}
The denominators $q_n$ tend to infinity. Given $M\geq1$, we have
$A_n=n+3\geq M$ for all sufficiently large $n$, and hence
\[
0<\abs{\xi-r_n}<q_n^{-M}
\]
for infinitely many $n$. Thus $\xi\in\Lio$.

By \textup{(F3)}, the positive integers $T_n$ are strictly increasing
and tend to infinity. The half-open blocks
$[T_n,T_{n+1})\cap\Z_{>0}$ therefore partition all integers $b\geq T_0$.
For such a $b$, let $n$ be the unique index with
$T_n\leq b<T_{n+1}$. Since $\xi\in I_{n+1}$, property \textup{(F6)}
gives, for every $a\in\Z$,
\begin{equation}\label{eq:image-lower-en}
\abs{f(\xi)-a/b}>b^{-\tau}=b^{-100}.
\end{equation}
This is the asserted eventual lower bound.

For every $\lambda>\tau$, the inequality
$0<|f(\xi)-a/b|<b^{-\lambda}$ can hold only for $b<T_0$,
and for each such $b$ there are only finitely many possible
numerators. Hence $\mu(f(\xi))\leq\tau=100$. Finally, $\xi\in I_0\subseteq I_*\subseteq V$.
\qed

\begin{remark}\label{rem:escape-set}
The escape set
\[
\mathcal E_\tau(f)
:=
\{\xi\in\Lio\cap U:\mu(f(\xi))\leq\tau\}
\]
contains a Cantor subset in every nonempty open subinterval
of $U$. Indeed, the component $(u_{n+1},v_{n+1})$ used in the
inductive step has length at least $2R_n/\Lambda_{n+1}$.
Choose inside it two disjoint compact intervals, each of length
$R_n/(2\Lambda_{n+1})$. Both inherit the approximation and
exclusion estimates, and \textup{(F7)} holds for both after replacing
the factor $3$ in \textup{(Q4)} by $6$. Repeating this choice at every
node gives a nested binary construction whose interval
diameters tend to zero. Every branch yields a point of the
escape set, and the resulting set is a Cantor set.
\end{remark}

\subsection{Proof of Corollary~\ref{cor:intro-rigidity}}


If the conclusion were false, then $f$ would not be the restriction to
$U$ of any rational function in $\R(x)$. Theorem~\ref{thm:intro-local}
would then produce a point $\xi\in\Lio\cap U$ for which
$\mu(f(\xi))\leq100$. Hence $f(\xi)\notin\Lio$, contradicting
$f(\Lio\cap U)\subseteq\Lio$.
\qed


\subsection{Proof of Theorem~\ref{thm:intro-mahler}}


Let $F:\C\to\C$ be entire and suppose that
$F(\Lio)\subseteq\Lio$. Since $\Lio$ is dense in $\R$, every
$x\in\R$ is the limit of a sequence $(\xi_n)$ in $\Lio$. The values
$F(\xi_n)$ are real, and continuity gives
\[
F(x)=\lim_{n\to\infty}F(\xi_n)\in\R.
\]
Thus $F(\R)\subseteq\R$.

Apply Corollary~\ref{cor:intro-rigidity} to the real-analytic function
$F|_{\R}$. There are polynomials $P,Q\in\R[z]$, with $Q\neq0$ and with no
common zero in $\C$, such that
\[
F(x)=\frac{P(x)}{Q(x)}
\qquad(x\in\R).
\]
The entire function
\[
z\longmapsto Q(z)F(z)-P(z)
\]
vanishes on $\R$ and hence vanishes identically on $\C$. If $Q$ were
nonconstant, let $\zeta\in\C$ be a zero of $Q$. The identity
$QF-P\equiv0$ would give $P(\zeta)=0$, contrary to the coprimality of
$P$ and $Q$. Hence $Q$ is constant, and therefore $F\in\R[z]$.
\qed

\section{Arithmetic rigidity and nonlinear sharpness}
\label{sec:arithmetic-boundary}

Corollary~\ref{cor:intro-rigidity} leaves open the coefficient
field of rational functions with Maillet's property. We show
that a M\"obius transformation has this property exactly when
it belongs to $\mathbb Q(x)$, whereas in every higher degree
there are preserving polynomials with transcendental
coefficients.

\subsection{M\"obius rigidity}
\label{subsec:mobius-rigidity}

We prove Theorem~\ref{thm:intro-mobius}. Let
\begin{equation}\label{eq:mobius-map}
M(x)=\frac{ax+b}{cx+d},
\qquad
\Delta_M:=ad-bc\neq0,
\end{equation}
with real coefficients. Choose an index
$\ell\in\{1,2,3,4\}$ for which the corresponding coordinate of
\[
\mathbf u:=(a,b,c,d)
\]
is nonzero, and rescale \(\mathbf u\) so that its \(\ell\)-th coordinate
is equal to \(1\). We continue to denote the normalized coefficients by
\(a,b,c,d\), and put \(\Delta_M=ad-bc\).
Let \(\boldsymbol\theta\in\R^3\) be the vector of the remaining
coordinates, in their inherited order. Then
\begin{equation}\label{eq:mobius-rational-chart}
M\in\Q(x)
\quad\Longleftrightarrow\quad
\boldsymbol\theta\in\Q^3.
\end{equation}
For \(h\in\Z_{>0}\), define
\begin{equation}\label{eq:mobius-Delta}
\Delta_h(\boldsymbol\theta)
:=
\min_{\mathbf r\in\Z^3}
\left\|\boldsymbol\theta-\frac{\mathbf r}{h}\right\|_\infty.
\end{equation}
By \eqref{eq:mobius-rational-chart}, if \(M\notin\Q(x)\), then
\(\Delta_h(\boldsymbol\theta)>0\) for every \(h\geq1\).

The proof splits according to the approximation properties of this
coefficient vector. Arbitrarily strong rational approximations lead to
a Liouville point in a rational fiber. Otherwise, a Diophantine lower
bound, combined with three-point reconstruction, yields a two-height
counting estimate.

\begin{lemma}[Simultaneous approximation dichotomy]
\label{lem:mobius-dichotomy}
Assume that \(M\notin\Q(x)\). Exactly one of the following alternatives
holds:
\begin{enumerate}[label=\textup{(\alph*)},leftmargin=2.7em]
\item for every \(N,H\in\Z_{>0}\) there exists \(h\geq H\) such that
\[
\Delta_h(\boldsymbol\theta)<h^{-N};
\]
\item there exist an integer \(\nu\geq1\) and a constant \(\eta_0>0\)
such that
\begin{equation}\label{eq:mobius-barrier}
\left\|\boldsymbol\theta-\frac{\mathbf r}{h}\right\|_\infty
\geq \eta_0h^{-\nu}
\end{equation}
for every \(h\in\Z_{>0}\) and every \(\mathbf r\in\Z^3\).
\end{enumerate}
\end{lemma}

\begin{proof}
If \textup{(a)} fails, there exist $N,H\in\Z_{>0}$ such that
$\Delta_h(\boldsymbol\theta)\geq h^{-N}$ for every $h\geq H$.
Increasing $H$ if necessary, we may assume that $H\geq2$.
Since the finitely many numbers $\Delta_h(\boldsymbol\theta)$,
$1\leq h<H$, are positive, one may take $\nu=N$ and
\[
\eta_0:=
\min\left\{
1,\min_{1\leq h<H}h^N\Delta_h(\boldsymbol\theta)
\right\}>0.
\]
This gives \textup{(b)}. Conversely, \textup{(b)} plainly excludes
\textup{(a)}.
\end{proof}

In alternative \textup{(a)}, rational approximation of the coefficient
vector transfers to rational approximation of the inverse image of a
fixed rational value.

\begin{lemma}[Rational fibers in the projectively Liouville case]
\label{lem:mobius-rational-fiber}
Assume that \(M\notin\Q(x)\) and that alternative \textup{(a)} of
Lemma~\ref{lem:mobius-dichotomy} holds. Then every nonempty open
interval \(V\) on which \(M\) is defined contains a point \(\xi\in\Lio\)
such that
\[
M(\xi)\in\Q.
\]
\end{lemma}

\begin{proof}
The image \(M(V)\) is a nonempty open interval. Choose
\(\rho\in\Q\cap M(V)\) such that
\begin{equation}\label{eq:mobius-fiber-xi}
\xi:=M^{-1}(\rho)
\end{equation}
is irrational. Such a choice is possible. Indeed, if three distinct
rational values had rational preimages, then \(M\) would send three
distinct rational points to three distinct rational points. A M\"obius
transformation is uniquely determined by its values at three distinct
points, and the corresponding linear system has rational coefficients;
hence its projective coefficient vector would belong to
\(\mathbb P^3(\Q)\), contrary to
\eqref{eq:mobius-rational-chart}. By construction, \(\xi\in V\).

Using alternative \textup{(a)} recursively, choose strictly increasing
integers \(h_n\) and vectors \(\mathbf r_n\in\Z^3\) such that
\begin{equation}\label{eq:mobius-simultaneous-approximants}
\left\|\boldsymbol\theta-\frac{\mathbf r_n}{h_n}\right\|_\infty
<h_n^{-n}
\qquad(n\geq1).
\end{equation}
Insert the three coordinates of \(\mathbf r_n/h_n\) into the fixed
chart, keeping the \(\ell\)-th coordinate equal to \(1\), and write the
resulting coefficient vector as
\[
\mathbf u_n=(a^{(n)},b^{(n)},c^{(n)},d^{(n)}).
\]
Since \(\mathbf u_n\to\mathbf u\), we have
\(a^{(n)}d^{(n)}-b^{(n)}c^{(n)}\neq0\) for all sufficiently large \(n\).
Let \(M_n\in\Q(x)\) be the corresponding M\"obius transformation.
Since \(a-\rho c\neq0\), the map sending a coefficient vector near
\(\mathbf u\) to the finite preimage of \(\rho\) is smooth. Hence
\[
\xi_n:=M_n^{-1}(\rho)\in\Q,
\qquad
|\xi-\xi_n|\leq C_1h_n^{-n}
\]
for some \(C_1>0\) and all sufficiently large \(n\).

After multiplying \(\mathbf u_n\) by \(h_n\), all four coefficients are
integers of size \(O(h_n)\). Writing \(\rho=r/s\) in lowest terms and
using
\[
M_n^{-1}(\rho)=
\frac{\rho d^{(n)}-b^{(n)}}{a^{(n)}-\rho c^{(n)}},
\]
we see that the positive denominator \(q_n'\) of \(\xi_n\) in lowest
terms satisfies
\begin{equation}\label{eq:mobius-fiber-denominator}
q_n'\leq C_2h_n
\end{equation}
for some \(C_2>0\). Moreover, \(\xi_n\to\xi\), and \(\xi\) is
irrational, so \(q_n'\to\infty\).

Fix \(A\geq1\). For all sufficiently large \(n>A\),
\[
0<|\xi-\xi_n|
\leq C_1h_n^{-n}
<C_2^{-A}h_n^{-A}
\leq(q_n')^{-A}.
\]
Thus \(\xi\in\Lio\), while \(M(\xi)=\rho\in\Q\).
\end{proof}

We now turn to the Diophantine case. The following reconstruction
lemma replaces the determinant step in Section~\ref{sec:scarcity}:
three rational points sufficiently close to the graph of \(M\)
determine a rational M\"obius transformation with nearby coefficients
and controlled height.

\begin{lemma}[Three-point rational reconstruction]
\label{lem:mobius-reconstruction}
Fix the M\"obius transformation \(M\) in \eqref{eq:mobius-map}, with the
normalization above, and a nondegenerate compact interval \(I_*\) on
which \(M\) is defined. There exist constants \(c_0,C_0>0\), depending
only on \(M\) and \(I_*\), with the following property.

Let \(x_1,x_2,x_3\in I_*\) be distinct, let \(y_1,y_2,y_3\in\R\), and
put
\[
\varepsilon:=\max_{1\leq i\leq3}|y_i-M(x_i)|,
\qquad
V(x_1,x_2,x_3):=\prod_{1\leq i<j\leq3}|x_i-x_j|.
\]
If
\begin{equation}\label{eq:mobius-reconstruction-smallness}
\varepsilon\leq c_0V(x_1,x_2,x_3),
\end{equation}
then there is a coefficient vector \(\widehat{\mathbf u}\), normalized
in the same chart as \(\mathbf u\), such that
\begin{equation}\label{eq:mobius-reconstruction-distance}
\|\widehat{\mathbf u}-\mathbf u\|_\infty
\leq
C_0\frac{\varepsilon}{V(x_1,x_2,x_3)}.
\end{equation}
It represents a nonconstant M\"obius transformation
\(\widehat M\) satisfying \(\widehat M(x_i)=y_i\) for \(1\leq i\leq3\).

If, in addition, for some positive real numbers $Q,B$,
\[
x_i=\frac{p_i}{q_i},
\qquad
y_i=\frac{r_i}{b_i},
\qquad
q_i<2Q,
\qquad
b_i<2B,
\]
with \(p_i,r_i\in\Z\) and \(q_i,b_i\in\Z_{>0}\), then the three free
coordinates of \(\widehat{\mathbf u}\) can be written with a common
positive denominator \(h\) satisfying
\begin{equation}\label{eq:mobius-reconstruction-height}
1\leq h\leq C_0Q^3B^3.
\end{equation}
\end{lemma}

\begin{proof}
For \(x,y\in\R\), put
\[
v(x,y):=(x,1,-xy,-y)\in\R^4.
\]
Let \(\mathcal A_0\) be the \(3\times4\) matrix with rows
\(v(x_i,M(x_i))\), and let \(\mathcal A\) have rows \(v(x_i,y_i)\).
For a \(3\times4\) matrix \(\mathcal B\), let
\(w_j(\mathcal B):=(-1)^{j+1}\det\mathcal B^{(j)}\), where
\(\mathcal B^{(j)}\) is obtained by deleting column \(j\). Then
\(\mathcal B\mathbf w(\mathcal B)^{\mathsf T}=0\). A direct
calculation gives
\begin{equation}\label{eq:mobius-cofactor-formula}
\mathbf w(\mathcal A_0)
=\Lambda(x_1,x_2,x_3)\mathbf u,
\end{equation}
where
\[
\Lambda(x_1,x_2,x_3)
=
\frac{\Delta_M
(x_1-x_2)(x_1-x_3)(x_2-x_3)}
{(cx_1+d)(cx_2+d)(cx_3+d)}.
\]
Since \(I_*\) is compact and contains no pole of \(M\),
\begin{equation}\label{eq:mobius-Lambda-lower}
|\Lambda(x_1,x_2,x_3)|
\geq c_1V(x_1,x_2,x_3)
\end{equation}
for some \(c_1>0\).

First restrict \(c_0\leq1\). All entries of \(\mathcal A_0\) and,
under \eqref{eq:mobius-reconstruction-smallness}, of \(\mathcal A\)
then lie in a fixed compact set. The maximal minors, being polynomial
in the matrix entries, satisfy
\begin{equation}\label{eq:mobius-cofactor-perturbation}
\|\mathbf w(\mathcal A)-\mathbf w(\mathcal A_0)\|_\infty
\leq C_1\varepsilon,
\end{equation}
with \(C_1\) independent of any subsequent decrease in \(c_0\).
Choose \(c_0\) so small that the right-hand side is at most
\(|\Lambda|/2\). Since \(u_\ell=1\), we have
\(w_\ell(\mathcal A_0)=\Lambda\), and hence
\(|w_\ell(\mathcal A)|\geq|\Lambda|/2>0\). Define
\[
\widehat{\mathbf u}:=
\frac{\mathbf w(\mathcal A)}{w_\ell(\mathcal A)}.
\]
Equations \eqref{eq:mobius-cofactor-formula}--\eqref{eq:mobius-cofactor-perturbation}
give
\[
\|\widehat{\mathbf u}-\mathbf u\|_\infty
\leq
\frac{2C_1(1+\|\mathbf u\|_\infty)\varepsilon}{|\Lambda|},
\]
which proves \eqref{eq:mobius-reconstruction-distance}.
Since \(ad-bc\neq0\) is an open condition on the coefficient vector,
decreasing \(c_0\) if necessary ensures that \(\widehat{\mathbf u}\)
represents a nonconstant M\"obius transformation.
Its denominator cannot vanish at any \(x_i\), since the kernel
relation would then force its numerator to vanish there as well,
contradicting its nonzero determinant. Thus the kernel relation gives
\(\widehat M(x_i)=y_i\).

For the height assertion, multiply the \(i\)-th row of \(\mathcal A\)
by \(q_ib_i\). The resulting integer row is
\[
(p_ib_i,q_ib_i,-p_ir_i,-q_ir_i).
\]
The numbers \(x_i\) lie in \(I_*\), and
\eqref{eq:mobius-reconstruction-smallness} bounds the \(y_i\)
uniformly in terms of \(M\) and \(I_*\). Thus every entry of this
integer matrix is \(O_{M,I_*}(QB)\). Its cofactor vector
\(\mathbf k\in\Z^4\) spans the same one-dimensional kernel as
\(\mathbf w(\mathcal A)\), and each of its coordinates is
\(O_{M,I_*}(Q^3B^3)\). Since \(k_\ell\neq0\), the normalized
coordinates are \(k_j/k_\ell\); taking \(h=|k_\ell|\) proves
\eqref{eq:mobius-reconstruction-height}, after enlarging \(C_0\) if
necessary. No reduction of this common denominator is needed.
\end{proof}

For the Diophantine branch, assume that alternative \textup{(b)} of
Lemma~\ref{lem:mobius-dichotomy} holds, and fix
\(\nu,\eta_0\) as in \eqref{eq:mobius-barrier} for this \(M\).
Fix a nondegenerate compact interval \(I_*\) on which \(M\) is defined.
Let \(m_*\) and \(M_*\) be the minimum and maximum of \(|M'|\) on
\(I_*\), so that
\begin{equation}\label{eq:mobius-derivative-bounds}
0<m_*\leq|M'(x)|\leq M_*
\qquad(x\in I_*).
\end{equation}
Choose an integer
\begin{equation}\label{eq:mobius-tau-choice}
\tau_1\geq9\nu+18
\end{equation}
and, for integers \(A\geq3\) and \(Q\geq2\), put
\begin{equation}\label{eq:mobius-T}
T_1(Q,A):=Q^{A/(\tau_1-3)}.
\end{equation}
For a nondegenerate compact interval \(J\subseteq I_*\) and integers
\(H,Q\geq2\), define
\begin{equation}\label{eq:mobius-dangerous-set}
\begin{split}
\mathcal D_M(J;Q,A,H)
:=
\Bigl\{\frac pq\in J:\;&\gcd(p,q)=1,\quad Q\leq q<2Q,\\
&\exists\,r\in\Z,\ b\in\Z_{>0}
\text{ with }H\leq b<T_1(Q,A),\\
&\abs{M(p/q)-r/b}
\leq2b^{-\tau_1}+4M_*Q^{-A}
\Bigr\}.
\end{split}
\end{equation}
As in Section~\ref{sec:scarcity}, the source fractions are reduced,
whereas the target fractions need not be.

\begin{proposition}[M\"obius two-height counting estimate]
\label{prop:mobius-scarcity}
There exists a constant \(C_{\mathrm{mob}}>0\), depending only on
\(M\), \(I_*\), and \(\tau_1\), with the following property. For every
integer \(A\geq3\), there exist constants \(C_A>0\) and \(Q_A\in\Z_{>0}\)
such that
\begin{equation}\label{eq:mobius-scarcity}
\#\mathcal D_M(J;Q,A,H)
\leq
C_{\mathrm{mob}}Q^2H^{2-\tau_1}+C_AQ
\end{equation}
for every nondegenerate compact interval \(J\subseteq I_*\), every
integer \(H\geq2\), and every integer \(Q\geq Q_A\). The constants
\(C_A\) and \(Q_A\) are uniform in \(J\) and \(H\).
\end{proposition}

\begin{proof}
Fix \(A\geq3\). All implied constants below may depend on the fixed
data \(M,I_*,\nu,\eta_0,\tau_1\), but not on \(J,H,Q\) or the target
block; any additional dependence on \(A\) is indicated.
If \(T_1(Q,A)\leq H\), then
\[
\mathcal D_M(J;Q,A,H)=\varnothing.
\]
Otherwise, decompose \([H,T_1(Q,A))\cap\Z_{>0}\) into the dyadic blocks
\[
\mathcal B_k
:=[B_k,\min\{2B_k,T_1(Q,A)\})\cap\Z_{>0},
\qquad B_k:=2^kH,
\]
with \(B_k<T_1(Q,A)\), and let \(\mathcal D_k\) be the source
centers admitting a witness with denominator in \(\mathcal B_k\).
Then \(\#\mathcal D_M(J;Q,A,H)\leq\sum_k\#\mathcal D_k\).
We call a block small if \(B_k<Q^{1/2}\) and large otherwise.

\smallskip
\noindent\textit{Small target heights.}
For a small block with left endpoint \(B\), the approximation condition
and the boundedness of \(M\) on \(I_*\) show that there are \(O(B^2)\)
pairs \((r,b)\) that can occur as witnesses. For each such pair,
monotonicity and \eqref{eq:mobius-derivative-bounds} show that
\[
E_{r,b}
:=
\set{x\in J:
\abs{M(x)-r/b}\leq2b^{-\tau_1}+4M_*Q^{-A}}
\]
is empty, a singleton, or an interval of length
\[
|E_{r,b}|\ll b^{-\tau_1}+Q^{-A}.
\]
Consequently, their union \(E_B\) has \(O(B^2)\) components and
\[
|E_B|\ll B^{2-\tau_1}+B^2Q^{-A}.
\]
Lemma~\ref{lem:rational-union} gives
\[
\#\mathcal D_k
\ll
Q^2B^{2-\tau_1}+Q^{2-A}B^2+B^2.
\]
Summing over \(B=2^kH<Q^{1/2}\), we obtain
\begin{equation}\label{eq:mobius-small-total}
\sum_{B_k<Q^{1/2}}\#\mathcal D_k
\leq
C_{\mathrm{mob}}Q^2H^{2-\tau_1}+O_A(Q),
\end{equation}
with \(C_{\mathrm{mob}}\) independent of \(J,A,H,Q\).

\smallskip
\noindent\textit{Large target heights.}
Fix a large block with left endpoint \(B=Q^u\). Then
\begin{equation}\label{eq:mobius-u-range}
u\geq\frac12,
\qquad
u<\frac{A}{\tau_1-3}.
\end{equation}
We shall show that, for all sufficiently large \(Q\), this block
contains at most two resonant centers. Indeed, three such centers
would determine, via Lemma~\ref{lem:mobius-reconstruction}, a rational
M\"obius transformation sufficiently close to \(M\) to contradict the
Diophantine barrier \eqref{eq:mobius-barrier}. Suppose therefore that
the block contains three distinct resonant centers
\[
x_i=\frac{p_i}{q_i},
\qquad Q\leq q_i<2Q,
\qquad1\leq i\leq3,
\]
with witnesses
\[
y_i=\frac{r_i}{b_i},
\qquad B\leq b_i<2B.
\]
Put
\[
\delta:=2B^{-\tau_1}+4M_*Q^{-A},
\qquad
s:=\min\{\tau_1u,A\}.
\]
Then \(|y_i-M(x_i)|\leq\delta\ll Q^{-s}\), and
\begin{equation}\label{eq:mobius-s-lower}
s\geq(\tau_1-3)u.
\end{equation}
Since the source fractions are distinct and reduced,
\[
V(x_1,x_2,x_3)
=\prod_{i<j}|x_i-x_j|
>\frac{1}{64Q^6}.
\]
The choice \eqref{eq:mobius-tau-choice} implies
\(Q^6\delta\to0\), uniformly in the block. Hence
Lemma~\ref{lem:mobius-reconstruction} applies for all sufficiently
large \(Q\), and produces a rational coefficient vector in the fixed
chart satisfying
\begin{equation}\label{eq:mobius-coefficient-upper}
\|\widehat{\mathbf u}-\mathbf u\|_\infty
\ll Q^6\delta
\ll Q^{6-s}.
\end{equation}
Its free coordinates have a common positive denominator \(h\) with
\[
1\leq h\ll Q^3B^3=Q^{3(1+u)}.
\]
Applying \eqref{eq:mobius-barrier} to these coordinates gives
\begin{equation}\label{eq:mobius-coefficient-lower}
\|\widehat{\mathbf u}-\mathbf u\|_\infty
\gg Q^{-3\nu(1+u)}.
\end{equation}
The exponent gap between \eqref{eq:mobius-coefficient-upper} and
\eqref{eq:mobius-coefficient-lower} satisfies
\begin{align}
 s-6-3\nu(1+u)
&\geq
(\tau_1-3-3\nu)u-(3\nu+6)\notag\\
&\geq
\frac{\tau_1-3-3\nu}{2}-(3\nu+6)
\geq\frac32,
\label{eq:mobius-exponent-gap}
\end{align}
where we used \(u\geq1/2\) and \eqref{eq:mobius-tau-choice}.
Thus the reconstructed coefficient vector would violate the
Diophantine barrier for all sufficiently large \(Q\), uniformly in the
block, in \(J\), and in \(H\). Each large block therefore contains at
most two resonant centers; no subdivision into short source cells is
needed.

There are at most
\[
1+\log_2\frac{T_1(Q,A)}H
\leq1+\frac{A}{\tau_1-3}\log_2Q
\]
large blocks, so their total contribution is \(O_A(\log Q)=O_A(Q)\).
Together with \eqref{eq:mobius-small-total}, this proves
\eqref{eq:mobius-scarcity}. Every lower threshold imposed on \(Q\)
depends only on the fixed data and \(A\), not on \(J\) or \(H\).
\end{proof}

The preceding counting estimate supplies safe rational centers at every
sufficiently large source scale. A nested-interval construction now
turns these centers into the required Liouville point.

\begin{lemma}[Finite-exponent escape in the Diophantine case]
\label{lem:mobius-finite-escape}
Assume that alternative \textup{(b)} of
Lemma~\ref{lem:mobius-dichotomy} holds. Then every nonempty open
interval \(V\) on which \(M\) is defined contains a point
\(\xi\in\Lio\) such that
\[
\mu(M(\xi))\leq\tau_1.
\]
\end{lemma}

\begin{proof}
We follow the fusion argument of Section~\ref{sec:fusion}, with
Proposition~\ref{prop:mobius-scarcity} supplying the counting estimate.
There are no Wronskian zero sets to avoid, so the choice of the next
interval is simpler.

Choose a nondegenerate compact interval \(I_0\Subset V\), and use it as
\(I_*\) in Proposition~\ref{prop:mobius-scarcity}. Write \(J(I)\) for
the middle third of \(I\), as in Section~\ref{sec:fusion}.
We construct nested nondegenerate compact intervals \(I_n\), strictly
increasing positive integers \(T_n\), and reduced fractions
\(r_n=p_n/q_n\).

Put \(A_n:=n+3\) and \(J_n:=J(I_n)\). Choose \(T_0\geq2\) so large
that
\begin{equation}\label{eq:mobius-initial-tail}
C_{\mathrm{mob}}T_0^{2-\tau_1}
\leq\frac{c_{\mathrm F}}4|J_0|.
\end{equation}
Suppose that \(I_n\) and \(T_n\) have been chosen and satisfy
\begin{equation}\label{eq:mobius-tail-invariant}
C_{\mathrm{mob}}T_n^{2-\tau_1}
\leq\frac{c_{\mathrm F}}4|J_n|.
\end{equation}
Proposition~\ref{prop:mobius-scarcity},
Lemma~\ref{lem:farey-lower-en}, and
\eqref{eq:mobius-tail-invariant} imply that every sufficiently large
\(Q\) admits a reduced fraction
\[
r=\frac pq\in J_n,
\qquad Q\leq q<2Q,
\]
satisfying
\begin{equation}\label{eq:mobius-safe-center}
\abs{M(r)-a/b}
>
2b^{-\tau_1}+4M_*Q^{-A_n}
\end{equation}
for every \(a\in\Z\) and every integer \(b\) with
\(T_n\leq b<T_1(Q,A_n)\).

Choose an integer \(Q_n\geq2\) large enough for this conclusion and,
simultaneously,
\begin{align}
Q_n&>2q_{n-1}\quad(n\geq1),
\label{eq:mobius-Q1}\\
(2Q_n)^{-A_n}&<\frac{|I_n|}{3},
\label{eq:mobius-Q2}\\
Q_n^{A_n/(\tau_1-3)}&\geq2T_n+2,
\label{eq:mobius-Q3}\\
12C_{\mathrm{mob}}2^{A_n+\tau_1-2}
Q_n^{-A_n/(\tau_1-3)}
&\leq\frac{c_{\mathrm F}}4.
\label{eq:mobius-Q4}
\end{align}
All quantities other than \(Q_n\) have already been fixed, and the
left-hand side of \eqref{eq:mobius-Q4} tends to zero. Select a safe
center
\[
r_n=\frac{p_n}{q_n}\in J_n,
\qquad Q_n\leq q_n<2Q_n,
\]
and put
\[
R_n:=(2Q_n)^{-A_n},
\qquad
T_{n+1}:=
\left\lfloor Q_n^{A_n/(\tau_1-3)}\right\rfloor.
\]
Condition \eqref{eq:mobius-Q3} gives \(T_{n+1}>T_n\).
Define
\begin{equation}\label{eq:mobius-next-interval}
I_{n+1}:=
\left[r_n+\frac{R_n}{4},r_n+\frac{R_n}{2}\right].
\end{equation}
Since \(r_n\in J_n\), condition \eqref{eq:mobius-Q2} gives
\(I_{n+1}\subset\operatorname{int}(I_n)\). For every \(x\in I_{n+1}\),
\begin{equation}\label{eq:mobius-source-approximation}
0<|x-r_n|<R_n<q_n^{-A_n}.
\end{equation}
Moreover, if \(T_n\leq b<T_{n+1}\), then
\(b<T_1(Q_n,A_n)\), and \eqref{eq:mobius-safe-center} together with
\(|M(x)-M(r_n)|\leq M_*|x-r_n|\) gives
\begin{equation}\label{eq:mobius-target-avoidance}
\abs{M(x)-a/b}>b^{-\tau_1}
\qquad(a\in\Z,\ x\in I_{n+1}).
\end{equation}

The interval \(I_{n+1}\) has length \(R_n/4\), so
\(|J_{n+1}|=R_n/12\). Since
\(T_{n+1}\geq\frac12Q_n^{A_n/(\tau_1-3)}\), condition
\eqref{eq:mobius-Q4} gives
\[
\frac{C_{\mathrm{mob}}T_{n+1}^{2-\tau_1}}{|J_{n+1}|}
\leq
12C_{\mathrm{mob}}2^{A_n+\tau_1-2}
Q_n^{-A_n/(\tau_1-3)}
\leq\frac{c_{\mathrm F}}4.
\]
This is \eqref{eq:mobius-tail-invariant} at stage \(n+1\), completing
the inductive step.

The nested intervals have diameters tending to zero, and hence
\(\bigcap_{n\geq0}I_n=\{\xi\}\) for some \(\xi\in V\). From
\eqref{eq:mobius-Q1} and \eqref{eq:mobius-source-approximation},
\(q_n\to\infty\) and
\[
0<|\xi-r_n|<q_n^{-A_n}.
\]
Since \(A_n\to\infty\), we have \(\xi\in\Lio\). Finally, the blocks
\([T_n,T_{n+1})\cap\Z_{>0}\) cover all sufficiently large denominators, and
\eqref{eq:mobius-target-avoidance} gives
\[
\abs{M(\xi)-a/b}>b^{-\tau_1}
\qquad(a\in\Z,\ b\geq T_0).
\]
Therefore \(\mu(M(\xi))\leq\tau_1\).
\end{proof}

Together with Lemma~\ref{lem:mobius-rational-fiber}, this treats both
alternatives in Lemma~\ref{lem:mobius-dichotomy} and yields the
degree-one classification.

\begin{proof}[Proof of Theorem~\ref{thm:intro-mobius}]
If \(M\in\Q(x)\), then Maillet's theorem gives Maillet's property.
Conversely, suppose that \(M\notin\Q(x)\). By
Lemma~\ref{lem:mobius-dichotomy}, either alternative \textup{(a)} or
alternative \textup{(b)} holds. In the first case,
Lemma~\ref{lem:mobius-rational-fiber} produces, in every open interval
on which \(M\) is defined, a Liouville number whose image is rational.
In the second case, Lemma~\ref{lem:mobius-finite-escape} produces a
Liouville number whose image has finite irrationality exponent. In
either case Maillet's property fails. This proves both assertions.
\end{proof}

\subsection{Nonlinear sharpness}
\label{subsec:nonlinear-sharpness}

M\"obius rigidity cannot be extended to higher degree. Recall that a
Liouville number \(\sigma\), with continued-fraction convergents
\(p_n/q_n\), is called \emph{strong Liouville} if
\[
\abs{\sigma-p_n/q_n}<q_n^{-\omega_n}
\]
for some sequence \(\omega_n\to\infty\).
Equivalently (see \cite{Petruska1992}),
\begin{equation}\label{eq:strong-convergent-criterion}
\frac{\log q_{n+1}}{\log q_n}\longrightarrow\infty.
\end{equation}

We shall use the following consequence of Petruska's theorem
\cite[Theorem~1]{Petruska1992}: if $\sigma$ is strong Liouville
and $y\in\Lio$, then
\begin{equation}\label{eq:Petruska-product}
\sigma y\in\Q\cup\Lio.
\end{equation}
The rational alternative in \eqref{eq:Petruska-product} is essential.
To exclude it for the monomials constructed below, we shall place their
nonzero rational fibers in Mahler's classes \(U_m\), with \(m\geq2\).

Here $\mathrm h(\alpha)$ denotes the \emph{absolute logarithmic Weil
height} of an algebraic number $\alpha$. We use the following
sequential height formulation: a transcendental number $\kappa$
belongs to $U_m$ if $m$ is the least positive integer for which there
exist algebraic numbers $\beta_n$ of degree $m$ and positive real
numbers $w_n$ such that
\[
\mathrm h(\beta_n)\to\infty,\qquad
w_n\to\infty,\qquad
|\kappa-\beta_n|
\leq e^{-w_n\mathrm h(\beta_n)}.
\]

The preceding formulation is the absolute-logarithmic-height version
of LeVeque's definition of $U_m$-numbers
\cite[Remark~2.6]{BiluMarquesMoreira2026} and, by the equivalence of
Mahler's and Koksma's classifications
\cite[Theorem~3.6]{Bugeaud2004}, defines the usual $U_m$-classes. In
particular, $U_1=\Lio$, and the classes $U_m$ are pairwise disjoint. We shall use the following gap principle in the form of
\cite[Theorem~3.1 and Corollary~3.3]{BiluMarques2026}.

\begin{lemma}[$U_m$ gap principle]
\label{lem:Um-gap-principle}
Let \(m\geq1\), let \(\kappa\in\C\) be transcendental, and let
\((\beta_n)\) be algebraic numbers of degree at most \(m\). Put
\(H_n:=\mathrm h(\beta_n)\), and let \((w_n)\) be a sequence of
positive real numbers. Suppose that
\[
H_n\to\infty,
\qquad
w_n\to\infty,
\qquad
|\kappa-\beta_n|\leq Ce^{-w_nH_n},
\qquad
H_{n+1}\leq Bw_nH_n
\]
for fixed constants \(B,C>0\). If \(\beta_n\) has degree exactly
\(m\) for all sufficiently large \(n\), then \(\kappa\in U_m\).
\end{lemma}

We now construct strong Liouville numbers for which the preceding gap
principle applies simultaneously to every nonzero rational fiber and
every degree \(m\geq2\). The denominator congruences enforce
irreducibility in every degree, while rapid growth ensures strong
Liouville approximation.

\begin{theorem}[Nonlinear sharpness]
\label{thm:strong-nonlinear-preservers}
There exist continuum many positive strong Liouville numbers \(\sigma\)
with the following property: for every integer \(m\geq2\) and every
\(r\in\Q^\times\), each root of
\begin{equation}\label{eq:nonlinear-fiber-equation}
X^m-\frac r\sigma
\end{equation}
is a \(U_m\)-number. Consequently, for every such \(\sigma\),
\begin{equation}\label{eq:nonlinear-preservation}
\sigma\xi^m\in\Lio
\qquad(\xi\in\Lio,\ m\geq2).
\end{equation}
In particular, in every degree \(m\geq2\) there are
Liouville-preserving polynomials in
\(\R[X]\setminus\overline{\Q}[X]\).
\end{theorem}

\begin{proof}
Enumerate
\[
\Q^\times=\set{r_1,r_2,\ldots},
\qquad
r_j=\frac{A_j}{B_j},
\]
where \(A_j\in\Z\setminus\{0\}\), \(B_j\in\Z_{>0}\), and
\(\gcd(A_j,B_j)=1\). We construct
\[
\sigma=[1;a_1,a_2,\ldots],
\qquad
\frac{p_n}{q_n}=[1;a_1,\ldots,a_n],
\]
with \(q_{-1}=0\) and \(q_0=1\). Suppose that
\(a_1,\ldots,a_{n-1}\) have been chosen. Select a prime
\(\ell_n\) such that
\begin{equation}\label{eq:ell-avoidance}
\ell_n\nmid q_{n-1}\prod_{j=1}^nA_jB_j.
\end{equation}
Because \(q_{n-1}\) is invertible modulo \(\ell_n^2\), one may choose a
positive integer \(a_n\), arbitrarily large, such that
\begin{equation}\label{eq:ell-valuation}
q_n=a_nq_{n-1}+q_{n-2}
\equiv\ell_n\pmod{\ell_n^2}
\end{equation}
and
\begin{equation}\label{eq:q-strong-growth}
q_n>q_{n-1}^{n+1}.
\end{equation}
Thus \(\ord_{\ell_n}(q_n)=1\), while
\(\ell_n\nmid p_n\) because \(\gcd(p_n,q_n)=1\). Condition
\eqref{eq:q-strong-growth} implies
\eqref{eq:strong-convergent-criterion}; hence \(\sigma\) is strong
Liouville, and \(1<\sigma<2\).

Fix \(j\geq1\), \(m\geq2\), and write
\(\sigma_n=p_n/q_n\). For \(n\geq j\), the polynomial
\begin{equation}\label{eq:Eisenstein-polynomial}
G_{j,m,n}(X):=B_jp_nX^m-A_jq_n
\end{equation}
is Eisenstein at \(\ell_n\). Indeed, \(\ell_n\) divides the constant
coefficient exactly once and does not divide the leading coefficient;
the intermediate coefficients vanish. Hence every root of
\[
X^m-\frac{r_j}{\sigma_n}
\]
has degree exactly \(m\) over \(\Q\).

Fix a root \(\beta\in\C\) satisfying \(\beta^m=r_j/\sigma\).
Since \(\sigma\) is transcendental, so is \(\beta\).
The holomorphic implicit-function theorem, applied to
\(zw^m-r_j=0\) at \((\sigma,\beta)\), supplies a local branch
\(u(z)\) with \(u(\sigma)=\beta\) and \(u(z)^m=r_j/z\). For all
sufficiently large \(n\), put \(\beta_n:=u(\sigma_n)\). Then
\([\Q(\beta_n):\Q]=m\), and, since \(1<\sigma,\sigma_n<2\),
\begin{equation}\label{eq:beta-approximation}
|\beta-\beta_n|
\ll_{j,m}\abs{\sigma-\sigma_n}
<\frac{C_{j,m}}{q_nq_{n+1}}.
\end{equation}
Moreover,
\[
m\,\mathrm h(\beta_n)
=
\mathrm h\left(\frac{A_jq_n}{B_jp_n}\right).
\]
The greatest common divisor of \(A_jq_n\) and \(B_jp_n\) divides
\(|A_jB_j|\), while \(1<p_n/q_n<2\). Therefore, with
\(H_n:=\mathrm h(\beta_n)\),
\begin{equation}\label{eq:beta-height-asymptotic}
H_n=\frac1m\log q_n+O_{j,m}(1).
\end{equation}
For sufficiently large \(n\), we have \(H_n>0\); set
\[
w_n:=\frac{\log q_{n+1}}{2H_n}.
\]
Then \(w_n\to\infty\), and \eqref{eq:beta-approximation} gives
\[
|\beta-\beta_n|\leq C_{j,m}e^{-w_nH_n}.
\]
Since \(m\geq2\), equation \eqref{eq:beta-height-asymptotic} also gives
\[
H_{n+1}\leq\log q_{n+1}=2w_nH_n
\]
for all sufficiently large \(n\). After discarding finitely many
terms, Lemma~\ref{lem:Um-gap-principle} therefore yields
\(\beta\in U_m\). Since \(j,m\), and the chosen root were arbitrary,
the asserted fiber property follows.

Now let \(\xi\in\Lio\) and \(m\geq2\). Maillet's theorem gives
\(\xi^m\in\Lio\), and \eqref{eq:Petruska-product} gives
\[
\sigma\xi^m\in\Q\cup\Lio.
\]
If \(\sigma\xi^m=r\in\Q\), then \(r\neq0\) and \(\xi\) is a root of
\(X^m-r/\sigma\); hence \(\xi\in U_m\), contrary to
\(\xi\in U_1\). Thus \eqref{eq:nonlinear-preservation} holds.

Finally, at every stage the prescribed congruence class modulo
\(\ell_n^2\) contains infinitely many positive integers satisfying
\eqref{eq:q-strong-growth}. Retaining two choices at each stage gives
a binary tree of continuum many continued fractions, every one of which
satisfies the preceding argument.
\end{proof}

\begin{proof}[Proof of Theorem~\ref{thm:intro-nonlinear}]
Apply Theorem~\ref{thm:strong-nonlinear-preservers} and retain
\eqref{eq:nonlinear-preservation}.
\end{proof}

Theorems~\ref{thm:intro-mobius} and
\ref{thm:strong-nonlinear-preservers} show that coefficient rigidity
has a sharp break between degree one and higher degree. They do not
classify all rational preservers, but they show that the real
coefficient field in Theorem~\ref{thm:intro-mahler} is unavoidable:
the conclusion \(F\in\R[z]\) cannot be strengthened to
\(F\in\overline{\Q}[z]\).

\section*{Acknowledgments}

%This problem has accompanied me for more than one decade, and I am grateful to the many colleagues with whom I have had the opportunity to discuss it over that period. Their questions, comments, and perspectives have been valuable at various stages of this work. 

The author would like to thank Yuri Bilu, Carlos Gustavo Moreira, Michel Waldschmidt, Johannes Schleischitz and Jean Lelis for generous and stimulating discussions on this problem over the last decade. This work was supported by the Brazilian National Council for Scientific and Technological Development (CNPq), Grant No.~304467/2023-5.


\section*{Use of artificial intelligence tools}

OpenAI's ChatGPT was used solely for English-language editing and stylistic corrections. The author assumes full responsibility for the manuscript.


\begin{thebibliography}{99}

\bibitem{AlniacikSaias1994}
K.~Alnia\c{c}{\i}k and E.~Saias,
Une remarque sur les $G_\delta$-denses,
\emph{Arch. Math. (Basel)} \textbf{62} (1994), no.~5, 425--426.

\bibitem{Apostol1976}
T.~M. Apostol,
\emph{Introduction to Analytic Number Theory},
Undergraduate Texts in Mathematics,
Springer-Verlag, New York--Heidelberg, 1976.

\bibitem{BeresnevichDickinsonVelani2007}
V.~Beresnevich, D.~Dickinson, and S.~Velani,
Diophantine approximation on planar curves and the distribution of rational points,
\emph{Ann. of Math. (2)} \textbf{166} (2007), no.~2, 367--426,
with an appendix by R.~C. Vaughan.

\bibitem{BiluMarques2026}
Y.~Bilu and D.~Marques,
Values of algebraic functions at Liouville numbers,
\emph{Bull. Braz. Math. Soc. (N.S.)} \textbf{57} (2026), article no.~29.

\bibitem{BiluMarquesMoreira2026}
Y.~Bilu, D.~Marques, and C.~G.~Moreira,
Values of generalized Liouville power series at algebraic numbers,
\emph{Funct. Approx. Comment. Math.} \textbf{75} (2026), no.~1,
95--110.

\bibitem{Bocher1900}
M.~B{\^o}cher,
The theory of linear dependence,
\emph{Ann. of Math. (2)} \textbf{2} (1900--1901), no.~1/4, 81--96.

\bibitem{BombieriPila1989}
E.~Bombieri and J.~Pila,
The number of integral points on arcs and ovals,
\emph{Duke Math. J.} \textbf{59} (1989), no.~2, 337--357.

\bibitem{BostanDumas2010}
A.~Bostan and P.~Dumas,
Wronskians and linear independence,
\emph{Amer. Math. Monthly} \textbf{117} (2010), no.~8, 722--727.

\bibitem{BoxallJones2015}
G.~J. Boxall and G.~O. Jones,
Algebraic values of certain analytic functions,
\emph{Int. Math. Res. Not. IMRN} 2015, no.~4, 1141--1158.

\bibitem{Bugeaud2004}
Y.~Bugeaud,
\emph{Approximation by Algebraic Numbers},
Cambridge Tracts in Mathematics, vol.~160,
Cambridge University Press, Cambridge, 2004.

\bibitem{ComteYomdin2018}
G.~Comte and Y.~Yomdin,
Zeroes and rational points of analytic functions,
\emph{Ann. Inst. Fourier (Grenoble)} \textbf{68} (2018), no.~6, 2445--2476.

%\bibitem{Erdos1962}
%P.~Erd\H{o}s,
%Representations of real numbers as sums and products of %Liouville numbers,
%\emph{Michigan Math. J.} \textbf{9} (1962), 59--60.

\bibitem{FischlerRivoal2024}
S.~Fischler and T.~Rivoal,
Values of $E$-functions are not Liouville numbers,
\emph{J. \'{E}c. polytech. Math.} \textbf{11} (2024), 1--18.

\bibitem{Habegger2018}
P.~Habegger,
Diophantine approximations on definable sets,
\emph{Selecta Math. (N.S.)} \textbf{24} (2018), no.~2, 1633--1675.

\bibitem{Huang2015}
J.-J.~Huang,
Rational points near planar curves and Diophantine approximation,
\emph{Adv. Math.} \textbf{274} (2015), 490--515.

\bibitem{Huxley1994}
M.~N. Huxley,
The rational points close to a curve,
\emph{Ann. Scuola Norm. Sup. Pisa Cl. Sci. (4)} \textbf{21} (1994), no.~3, 357--375.

\bibitem{KleinbockMargulis1998}
D.~Y. Kleinbock and G.~A. Margulis,
Flows on homogeneous spaces and Diophantine approximation on manifolds,
\emph{Ann. of Math. (2)} \textbf{148} (1998), no.~1, 339--360.

%\bibitem{LelisMarques2020}
%J.~Lelis and D.~Marques,
%On transcendental entire functions mapping $\Q$ into itself,
%\emph{J. Number Theory} \textbf{206} (2020), 310--319.

\bibitem{LelisMoreiraSilva2026}
J.~Lelis, C.~G. Moreira, and E.~Silva,
On $C^k$-functions mapping $\Q$ into itself and Mahler's problem on Liouville numbers,
preprint, arXiv:2607.24427, 2026.

\bibitem{Liouville1844}
J.~Liouville,
Communication du 13 mai 1844,
\emph{C. R. Acad. Sci.} \textbf{18} (1844), 883--885, 910--911.

\bibitem{Liouville1851}
J.~Liouville,
Sur des classes tr\`es-\'etendues de quantit\'es dont la valeur
n'est ni alg\'ebrique, ni m\^eme r\'eductible \`a des
irrationnelles alg\'ebriques,
\emph{J. Math. Pures Appl. (1)} \textbf{16} (1851), 133--142.

\bibitem{Lombardo2017}
D.~Lombardo,
On the analytic bijections of the rationals in $[0,1]$,
\emph{Atti Accad. Naz. Lincei Rend. Lincei Mat. Appl.} \textbf{28} (2017),
no.~1, 65--83.

\bibitem{Mahler1984}
K.~Mahler,
Some suggestions for further research,
\emph{Bull. Aust. Math. Soc.} \textbf{29} (1984), no.~1, 101--108.

\bibitem{Maillet1906}
E.~Maillet,
\emph{Introduction \`a la th\'eorie des nombres transcendants et des propri\'et\'es
arithm\'etiques des fonctions},
Gauthier--Villars, Paris, 1906.

\bibitem{MarquesMoreira2015}
D.~Marques and C.~G. Moreira,
On a variant of a question proposed by K.~Mahler concerning Liouville numbers,
\emph{Bull. Aust. Math. Soc.} \textbf{91} (2015), no.~1, 29--33.

\bibitem{MarquesSchleischitz2016}
D.~Marques and J.~Schleischitz,
On a problem posed by Mahler,
\emph{J. Aust. Math. Soc.} \textbf{100} (2016), no.~1, 86--107.

\bibitem{MarquesSmooth2026}
D.~Marques,
Smooth diffeomorphisms and Mahler's problem on Liouville numbers,
preprint, arXiv:2609.00376, 2026.

\bibitem{Petruska1992}
G.~Petruska,
On strong Liouville numbers (a problem of P.~Erd\H{o}s),
\emph{Indag. Math. (N.S.)} \textbf{3} (1992), no.~2, 211--218.

\bibitem{Pila1991}
J.~Pila,
Geometric postulation of a smooth function and the number of rational points,
\emph{Duke Math. J.} \textbf{63} (1991), no.~2, 449--463.

\bibitem{Pila2005}
J.~Pila,
Rational points on a subanalytic surface,
\emph{Ann. Inst. Fourier (Grenoble)} \textbf{55} (2005), no.~5, 1501--1516.

\bibitem{PilaWilkie2006}
J.~Pila and A.~J. Wilkie,
The rational points of a definable set,
\emph{Duke Math. J.} \textbf{133} (2006), no.~3, 591--616.

%\bibitem{Schleischitz2025}
%J.~Schleischitz,
%Maillet's property and Mahler's conjecture on Liouville numbers fail for matrices,
%\emph{J. Aust. Math. Soc.} \textbf{119} (2025), no.~1, 106--127.

\bibitem{Stackel1895}
P.~St\"ackel,
Ueber arithmetische Eigenschaften analytischer Functionen,
\emph{Math. Ann.} \textbf{46} (1895), 513--520.

\bibitem{Stackel1902}
P.~St\"ackel,
Arithmetische Eigenschaften analytischer Functionen,
\emph{Acta Math.} \textbf{25} (1902), 371--383.

\bibitem{Surroca2006}
A.~Surroca,
Valeurs alg\'ebriques de fonctions transcendantes,
\emph{Int. Math. Res. Not.} 2006, Art. ID 16834, 31 pp.

\bibitem{VanderPoorten1968}
A.~J. van der Poorten,
Transcendental entire functions mapping every algebraic number field into itself,
\emph{J. Austral. Math. Soc.} \textbf{8} (1968), no.~2, 192--193.

\bibitem{VaughanVelani2006}
R.~C. Vaughan and S.~Velani,
Diophantine approximation on planar curves: the convergence theory,
\emph{Invent. Math.} \textbf{166} (2006), no.~1, 103--124.

\bibitem{Waldschmidt2003}
M.~Waldschmidt,
Algebraic values of analytic functions,
\emph{J. Comput. Appl. Math.} \textbf{160} (2003), no.~1--2, 323--333.

\bibitem{Waldschmidt2004}
M.~Waldschmidt,
Open Diophantine problems,
\emph{Mosc. Math. J.} \textbf{4} (2004), no.~1, 245--305.

\bibitem{Weierstrass1923}
K.~Weierstrass,
Briefe von K.~Weierstrass an L.~Koenigsberger,
\emph{Acta Math.} \textbf{39} (1923), 226--239.

\end{thebibliography}

\end{document}
